\documentclass[11pt]{article}

\usepackage[T1]{fontenc}
\usepackage{lmodern}
\usepackage{amsmath,amssymb,amsthm,mathtools}
\usepackage{mathrsfs}
\usepackage{booktabs}
\usepackage{array}
\usepackage{enumitem}
\usepackage[a4paper,margin=27mm]{geometry}
\usepackage[colorlinks=true,linkcolor=blue,citecolor=blue,urlcolor=blue]{hyperref}
\setlength{\emergencystretch}{2em}

\newtheorem{theorem}{Theorem}[section]
\newtheorem{proposition}[theorem]{Proposition}
\newtheorem{lemma}[theorem]{Lemma}
\newtheorem{corollary}[theorem]{Corollary}
\theoremstyle{definition}
\newtheorem{definition}[theorem]{Definition}
\newtheorem{construction}[theorem]{Construction}
\newtheorem{conjecture}[theorem]{Conjecture}
\newtheorem{hypothesis}[theorem]{Hypothesis}
\newtheorem{problem}[theorem]{Problem}
\newtheorem{example}[theorem]{Example}
\newtheorem{warning}[theorem]{Warning}
\theoremstyle{remark}
\newtheorem{remark}[theorem]{Remark}

\newcommand{\Qell}{\mathbf Q_\ell}
\newcommand{\Lmot}{\mathbb L_{\mathrm{mot}}}
\newcommand{\Lell}{\mathbb L_{\ell}}
\newcommand{\UZ}{U_{\mathbf Z}}
\newcommand{\Uk}{U_k}
\newcommand{\kappamot}{\kappa_{\mathrm{mot}}}
\newcommand{\kappaell}{\kappa_{\ell}}
\newcommand{\one}{\mathbf 1}
\newcommand{\cD}{\mathcal D}
\newcommand{\cP}{\mathcal P}
\newcommand{\cA}{\mathcal A}
\newcommand{\cK}{\mathcal K}
\newcommand{\cT}{\mathcal T}
\newcommand{\cX}{\mathcal X}
\newcommand{\cM}{\mathcal M}
\newcommand{\cF}{\mathcal F}
\newcommand{\cC}{\mathcal C}
\newcommand{\cR}{\mathcal R}
\newcommand{\D}{\mathbb D}
\newcommand{\Dcell}{\mathbb D_{\mathrm{cell}}}
\newcommand{\GF}{\operatorname{GF}}
\newcommand{\Cell}{\operatorname{Cell}}
\newcommand{\Trop}{\operatorname{Trop}}
\newcommand{\Sk}{\operatorname{Sk}}
\newcommand{\Spec}{\operatorname{Spec}}
\newcommand{\Cone}{\operatorname{Cone}}
\newcommand{\cofib}{\operatorname{cofib}}
\newcommand{\relint}{\operatorname{relint}}
\newcommand{\rk}{\operatorname{rk}}
\newcommand{\ord}{\operatorname{ord}}
\newcommand{\supp}{\operatorname{supp}}
\newcommand{\Tr}{\operatorname{Tr}}
\newcommand{\id}{\operatorname{id}}
\newcommand{\an}{\mathrm{an}}
\newcommand{\rt}{\mathrm{rt}}
\newcommand{\can}{\mathrm{can}}
\newcommand{\ad}{\mathrm{ad}}
\newcommand{\ab}{\mathrm{ab}}
\newcommand{\der}{\mathrm{der}}
\newcommand{\cO}{\mathcal O}
\newcommand{\vol}{\operatorname{vol}}
\newcommand{\Nm}{\operatorname{Nm}}


\title{From Represented $PGL_2$ Orbit Shells to Rank-One Orientation Reciprocity\\
\large Quadratic Residual Towers, Haar--Frobenius Comparison, and the
Unramified Absolute-$A_1$ Theorem}
\author{YINGFENG JIANG}
\date{28 August 2026}

\begin{document}
\maketitle

\begin{abstract}
Let
\[
 \UZ=\Spec\mathbf Z[1/2,a,a^{-1},(1-a)^{-1}],\qquad
 \mathcal E_{\mathbf Z}=\mathcal O_{\UZ}[z]/(z^2-a).
\]
For a field $k$ with $2,\ell\in k^\times$, write
$\Uk=\UZ\otimes_{\mathbf Z}k$ and
$\mathcal E_k=\mathcal E_{\mathbf Z}\otimes k$.  For every $r\geq1$ we
study the base change to $\Uk$ of the smooth affine truncated-unit quotient
\[
 \mathscr Q_r=
 \operatorname{Res}_{\mathcal E_k[\epsilon]/(\epsilon^r)/\Uk}\mathbb G_m
 \big/
 \operatorname{Res}_{\mathcal O_{\Uk}[\epsilon]/(\epsilon^r)/\Uk}\mathbb G_m.
\]
The first layer is the norm-one torus $\mathbb T$, and every subsequent
reduction is a torsor under the trace-zero vector line.  In a fixed enhanced
$\ell$-adic six-functor formalism, its compact cochains form the Frobenius
orientation bimodule
\[
 R\pi_{\mathbb T,!}\Qell
 \simeq R\pi_{\mathbb T,*}\Qell\otimes
 \det(X_*(\mathbb T)\otimes\Qell)[-1].
\]

Using the fixed lift
$\widetilde\gamma_m=1+\varpi^m\iota(z)\in
GL_2(\mathcal O_{\Uk}[[\varpi]])$
and its class $\gamma_m\in PGL_2$, we prove that $\mathscr Q_r$ represents
the compact positive-loop orbit sheaf of the quadratic order lattice
$\mathcal A_r=\mathcal O_{\Uk}[[\varpi]]+\varpi^r\mathcal E_k[[\varpi]]$ for
$0\leq r\leq m$.  This is an isomorphism of fppf sheaves over the entire
quadratic parameter space, not merely an equality of finite-field point
counts.  At field-valued points the quadratic-order classification gives the
full support decomposition of the spherical orbital integral.  The orbit
stabilizers, quotient Haar masses, and Frobenius traces agree term by term.

The split and unramified-elliptic fibers are thereby unified by the Kummer
local system
$\cK\simeq\det(X_*(\mathbb T)\otimes\Qell)$.  After Weyl normalization their
generating class satisfies
\[
 W_{PGL_2}(z)=\frac{1-\kappaell Qz}{(1-z)(1-Qz)},\qquad
 W_{PGL_2}(z^{-1})=-\kappaell z\,d_{\Uk}(W_{PGL_2}(z)).
\]
For every finite depth cutoff we also prove the object-level Verdier partner
identity on the discrete labelled-shell target.  The separate height-cone
statement is made only for the cohomological associated graded and after
passing to a rational $K_0$ character.

We then pass from the represented $PGL_2$ family to a second, explicitly
local level.  For central-neutral standard spherical families in unramified
absolute type $A_1$, a classification reduction and an orbit-summed measure
lemma exhaust the possible reduced-building shell lattices.  The $SL_2$
elliptic phases retain only even or odd spheres, whereas adjoining the
connected center in $U_2$ forces deletion of the central cocharacter
direction.  Nevertheless every case satisfies
\[
 W_G(z^{-1})=-z\,(\mathfrak o_{G,T}\otimes d_S)W_G(z),\qquad
 \mathfrak o_{G,T}
 =\det\!\left(X_*(T/Z(G)^\circ)\otimes\Qell\right).
\]
Outside the saturated $PGL_2$ family this is a local orbit-summed theorem,
not a relative fppf representability statement.  No whole affine-Springer
quotient, cross-shell incidence, convolution, Rees completion, or general
logarithmic skeletonization theorem is claimed.
\end{abstract}

\section{Introduction}

The numerical split and unramified-elliptic rank-one orbital integrals are
classical.  Likewise, the motivic nature of broad classes of orbital
integrals was established by Hales, and transfer principles place families
of fundamental-lemma integrals in motivic integration
\cite{HalesMotivic,CluckersHalesLoeser}.  Affine-Springer fundamental domains
and their relation to ordinary or weighted orbital integrals have been
developed systematically by Chen
\cite{ChenFundamentalDomain,ChenWeighted,ChenLocalBehavior}.  Accordingly,
neither the existence of a motivic interpolation nor the rational function
by itself is presented here as new.  Kummer local systems have also appeared
previously in orbital-integral geometry, notably in the Moy--Prasad quotient
analysis of Cunningham--Hales \cite{CunninghamHalesGood}.

The level isolated here is the representing-object level.
Hales constructs virtual motivic representatives for definable orbital
integrals, and the transfer principle controls identities after
specialization; neither identifies the individual compact $T(\mathcal O)$
orbits in this quadratic family with one fppf quotient over $U$.  Chen's
fundamental domains and truncations relate affine-Springer point counts to
ordinary or weighted orbital integrals, including explicit rank-two
calculations, but they do not provide the componentwise quotient
$\mathscr Q_r$ or its determinant-local-system descent.  Kottwitz's explicit
rank-one formulas supply a numerical consistency check
\cite{KottwitzOrbital}.  The theorem below proves the uniform represented shell
and the cohomological Kummer orientation, not a competing recalculation of
those classical specializations.

The first contribution is more specific and more geometric.  For the universal
quadratic $PGL_2$ family we represent each compact positive-loop orbit shell
by one smooth finite-type $U$-scheme $\mathscr Q_r$, prove the identification
as an fppf orbit sheaf, and track the split/elliptic variation through the
single determinant local system
\[
 \cK\simeq\det(X_*(\mathbb T)\otimes\Qell).
\]
The vector-torsor residual tower propagates this orientation coefficient
through every depth.  Thus the familiar specializations $q-1$ and $q+1$ are
fibers of one descended cohomological object rather than two separately
computed formulas.  The scope of the result is exactly this componentwise
represented-family theorem and its orientation-compatible assembly.

The second contribution is local and structural.  We classify the
reduced-building data of central-neutral standard spherical families for all
unramified connected reductive groups whose absolute adjoint root system is
$A_1$.  The classification has only a saturated shell lattice or a
type-preserving lattice with two elliptic embedding phases.  An orbit-summed
measure lemma handles the fact that one $PGL_2$ torus orbit may split into
several $SL_2$ orbits.  The $SL_2$ parity calculation changes the generating
series without changing its inversion line, while the $U_2$ center test shows
that the determinant must be taken after quotienting by $Z(G)^\circ$.  This
proves the unramified absolute-$A_1$ orientation theorem at the level of local
orbital masses and their displayed $K_0$ interpolation.

The logical spine of the paper is therefore
\[
 \boxed{
 \begin{gathered}
 \text{quadratic residual tower}
 \Longrightarrow
 \text{represented }PGL_2\text{ orbit shells}\\
 \Downarrow\\
 \text{Haar/Frobenius reciprocity}
 \Longrightarrow
 \text{finite shell skeletonization}
 \end{gathered}}
\]
The first implication is geometric, the second compares Haar mass with a
derived Frobenius supertrace, and the last is the fixed-map Verdier partner
for the finite labelled shell map.  Keeping these arrows distinct prevents
the rational height-cone shadow from being mistaken for an object-level
skeletonization of the whole affine-Springer fiber.

The structural extension is a second chain,
\[
 \begin{gathered}
 PGL_2\text{ orientation}
 \Longrightarrow
 \text{absolute-}A_1\text{ classification}\\
 \Downarrow\\
 SL_2\text{ parity and }U_2\text{ center test}
 \Longrightarrow
 \text{local rank-one orientation reciprocity}
 \end{gathered}
\]
Only the first chain contains a relative represented orbit family.

The paper is organized in the order needed to verify the two theorem layers.
Section~\ref{sec:main-theorem-status} fixes notation and states Theorems A
and B with their different logical status.  The enhanced torus calculation
and residual tower come next, followed by the $PGL_2$ orbit-shell comparison,
Haar support decomposition, and Frobenius trace.  The generating function,
finite labelled-shell Verdier partner, and cellular $K_0$ shadow are then
deduced.  Section~\ref{sec:absolute-A1} proves the classification reduction,
the $SL_2$ parity and $U_2$ center tests, and the local absolute-$A_1$
theorem.  The final section compares the result with the literature, states
the limitations, and isolates the remaining categorical/logarithmic bridge.

\subsection{Coefficient convention}

We fix, once and for all, the constructible $\Qell$-linear enhanced
six-functor formalism on the noetherian finite-dimensional schemes occurring
below, in the sense of the enhanced operations of
Laszlo--Olsson and Liu--Zheng \cite{LaszloOlsson,LiuZhengEnhanced}.
We use its symmetric monoidal stable categories $\mathscr D_c(S)$, internal
Hom, module categories for algebra objects, relative tensor products computed
by the bar construction, smooth purity, and effective finite \'etale descent.
Finite \'etale pullback is conservative; descent applies to $E_\infty$-algebras,
their module categories, and morphisms between dualizable modules.  These are
standing properties of the fixed formalism, not conclusions deduced from the
triangulated homotopy category.  We do not quantify over arbitrary
enhancements.

There are two Grothendieck levels, and we keep their notation separate.  Let
\[
 \widetilde\UZ=\Spec\mathbf Z[1/2,b,b^{-1},(1-b^2)^{-1}]
 \longrightarrow\UZ,\qquad b^2=a,
\]
and set
\[
 \Lmot=[\mathbb A^1_{\UZ}],\qquad
 \kappamot=[\widetilde\UZ]-[\UZ]
 \quad\text{in }K_0(\operatorname{Var}_{\UZ}).
\]
After base change to $\Uk$, compactly supported $\ell$-adic realization sends
these classes to
\[
 \Lell=[\Qell(-1)[-2]],\qquad
 \kappaell=[\cK]
 \quad\text{in }K_0(D_c^b(\Uk,\Qell)).
\]
Here $\cK$ is the anti-invariant summand defined by
$R\widetilde p_*\Qell\simeq\Qell\oplus\cK$ for the base-changed quadratic
cover.  The torsor identity gives $\kappamot^2=1$, hence
$\kappaell^2=1$ after realization.
We put $Q=\Lell^{-1}$.  Thus a geometric cover class, its Kummer-sheaf
realization, and a Tate class never share the same symbol below.

\subsection{Universal notation}

Fix $k$ as above.  From this point onward, until an integral model is
explicitly displayed, we abbreviate $U=\Uk$ and
$\mathcal E=\mathcal E_k$.  Let
$\iota:\mathcal E\hookrightarrow M_2(\mathcal O_U)$ be the regular
representation in the basis $(1,z)$, so
\[
 \iota(z)=\begin{pmatrix}0&a\\1&0\end{pmatrix}.
\]
Put $G=PGL_{2,U}$ and
\[
 \mathbb T=\operatorname{Res}_{\mathcal E/U}\mathbb G_m/\mathbb G_m
 \simeq\ker(N_{\mathcal E/\mathcal O_U}),\qquad
 \widetilde\gamma_m=1+\varpi^m\iota(z),\qquad
 \gamma_m=[\widetilde\gamma_m].
\]
Let $X_m\subset\operatorname{Gr}_G$ be the intrinsic torsor-valued
affine-Springer subfunctor defined in
\eqref{eq:PGL2-group-affine-Springer}; on a section represented by a loop
$g$, its condition is $g^{-1}\gamma_mg\in L^+G$.  Also put
$\mathcal A_r=\mathcal O_U[[\varpi]]+\varpi^r\mathcal E[[\varpi]]$.
All appearances of an order generated by $\gamma_m$ below mean the order
generated by the fixed $GL_2$-lift $\widetilde\gamma_m$.

The companion representative used in earlier calculations is
\[
 \gamma^{\mathrm{comp}}_{s,a}
 =\left[\begin{pmatrix}1&as^2\\1&1\end{pmatrix}\right].
\]
For $s=\varpi^m$ it is conjugate to the integral quadratic-basis model:
\begin{equation}
 \begin{pmatrix}1&0\\0&s\end{pmatrix}
 \begin{pmatrix}1&as^2\\1&1\end{pmatrix}
 \begin{pmatrix}1&0\\0&s^{-1}\end{pmatrix}
 =1+s\iota(z).
 \label{eq:companion-integral-conjugacy}
\end{equation}
The paper uses only $\gamma_m=[1+\varpi^m\iota(z)]$; the companion form is
recorded solely to identify the two presentations of the same regular class.

We reserve separate notation for the four orbital levels.  At a finite-field
point $x$, $\operatorname{OI}_m(x)$ denotes the numerical Haar integral.
The represented shell union will define the coefficient object
$\mathcal I_m^{\mathrm{orb}}=Rq^{\mathrm{orb}}_{m,!}\Qell$, its raw class is
$I_m^{\mathrm{orb}}=[\mathcal I_m^{\mathrm{orb}}]$, and its Weyl-normalized
class is
\[
 J_m^{\mathrm{orb}}=Q^mI_m^{\mathrm{orb}}.
\]
Only the normalized coefficients enter
$W_{PGL_2}(z)=\sum_{m\ge0}J_m^{\mathrm{orb}}z^m$.  Thus a numerical orbital
integral, a coefficient object, its Grothendieck class, and its normalization
are never denoted by the same symbol.

\begin{proposition}[The Kummer coefficient is strictly non-Tate]
\label{prop:strict-non-Tate-Kummer}
The local system $\cK$ on $\Uk$ does not belong to the thick tensor
subcategory generated by constant Tate sheaves.  Consequently every positive
residual layer contains a genuinely non-Tate cohomology sheaf; the correction
$\kappaell$ in the normalized orbital coefficients cannot be replaced by a
Tate power.
\end{proposition}

\begin{proof}
After base change to an algebraic closure, the cover $b^2=a$ of
$\mathbb G_m\setminus\{1\}$ is connected and has quadratic monodromy around
$a=0$.  Hence its anti-invariant summand $\cK$ has monodromy $-1$, whereas
constant Tate sheaves have trivial geometric monodromy.  The compact
cohomology calculation in Lemma~\ref{lem:torus-compact-cohomology}, followed
by the vector-torsor Thom isomorphism
\eqref{eq:Haar-shell-Thom}, places a Tate twist of $\cK$ in every
$Rq_{r,!}\Qell$ with $r\ge1$.  Weight and monodromy separate it from the Tate
summand, so it cannot disappear through a Tate re-expression.
\end{proof}

\begin{remark}
The non-Tate property of a nontrivial quadratic Kummer local system is
standard.  Proposition~\ref{prop:strict-non-Tate-Kummer} is included only to
certify that the coefficient occurring in the represented shells cannot be
absorbed into a Tate normalization.
\end{remark}


\section{Main results and status}
\label{sec:main-theorem-status}

There are two theorem layers.  Theorem A is relative and object-level for the
universal quadratic $PGL_2$ family.  Theorem B is a local orbit-summed theorem
for the explicitly defined unramified absolute-$A_1$ class.  We state both
before introducing the formal lemmas used later and keep their logical status
separate throughout.

\begin{theorem}[Theorem A: the universal quadratic $PGL_2$ package]
\label{thm:main-PGL2-package}
Let $k$ be a field with $2\in k^\times$, let $\ell$ be invertible in $k$, and
put
\[
 \Uk=\UZ\otimes_{\mathbf Z}k
 =\Spec k[a,a^{-1},(1-a)^{-1}],
 \qquad \mathcal E_k=\mathcal O_{\Uk}[z]/(z^2-a).
\]
For the universal family
$\gamma_m=[1+\varpi^m\iota(z)]\in
PGL_2(\mathcal O_{\Uk}((\varpi)))$ the following
statements hold.
\begin{enumerate}[label=\textup{(\roman*)}]
\item The truncated-unit fppf quotients $\mathscr Q_r$ are represented by
  smooth affine $U$-group schemes.  One has $\mathscr Q_1=\mathbb T$, the
  norm-one torus, and $\mathscr Q_r\to\mathscr Q_{r-1}$ is a torsor under the
  trace-zero vector line.  In the enhanced bimodule category,
  \[
   R\pi_{\mathbb T,!}\Qell
   \simeq R\pi_{\mathbb T,*}\Qell
   \otimes\det(X_*(\mathbb T)\otimes\Qell)[-1].
  \]
\item For every $0\le r\le m$, the positive-loop orbit sheaf of the order
  lattice
  $\mathcal A_r=\mathcal O_{\Uk}[[\varpi]]+\varpi^r\mathcal E_k[[\varpi]]$
  is represented, and its orbit morphism is a monomorphism of fppf sheaves
  \[
   \mathscr Q_r\xrightarrow{\sim}
   L^+\mathbb T\cdot[\mathcal A_r]\longrightarrow X_m.
  \]
  At every field-valued point, the classes $[A_r]$, $0\le r\le m$, are exactly
  the representatives of $T(F)\backslash X_m$.
\item With Haar measure normalized by
  $\operatorname{vol}(PGL_2(\mathcal O_F))=
   \operatorname{vol}(T(\mathcal O_F))=1$ and quotient measure defined by
  Weil integration, the mass of the $r$th double coset is
  \[
   [T(\mathcal O_F):A_r^\times/\mathcal O_F^\times].
  \]
  Over a finite residue field this is both $\#\mathscr Q_r(k)$ and the
  geometric-Frobenius derived supertrace of its compact coefficient.
\item After Weyl-discriminant normalization, the compact-orbital generating
  class is
  \begin{align*}
   J_m^{\mathrm{orb}}&=Q^mI_m^{\mathrm{orb}},\\
   W_{PGL_2}(z)&=\sum_{m\ge0}J_m^{\mathrm{orb}}z^m
   =\frac{1-\kappaell Qz}{(1-z)(1-Qz)},\\
   W_{PGL_2}(z^{-1})&=-\kappaell z\,d_{\Uk}(W_{PGL_2}(z)).
  \end{align*}
  where $I_m^{\mathrm{orb}}=[Rq^{\mathrm{orb}}_{m,!}\Qell]$ is the raw
  compact-orbital coefficient,
  $\kappaell=[\det(X_*(\mathbb T)\otimes\Qell)]=[\cK]$, and $d_{\Uk}$ is the
  dimension-normalized coefficient duality of
  \eqref{eq:normalized-U-duality}.
\item For every finite cutoff $N$, the actual labelled-shell map
  $\rho_N:\mathscr Y_N^{\mathrm{orb}}\to \Uk\times\Lambda_N$ defines
  $\Sk_{N,!}=R\rho_{N,!}$ and $\Sk_{N,*}=R\rho_{N,*}$ independently, and
  \[
   \D_{\mathrm{cell},N}\Sk_{N,!}
   \simeq
   \Sk_{N,*}\D_{\mathscr Y_N^{\mathrm{orb}}}
  \]
  naturally on the full bounded constructible category.
\item The associated-graded shell coefficients have canonical height
  $s=m-r+1$.  Their de-Tated height character is represented by the auxiliary
  cellular envelope
  $(1-\kappaell)\one_C+\kappaell\one_F$; the variable $u^s$ records the twist
  $(s)$, and the specialization $u\mapsto Q$ recovers the actual ordinary
  $K_0$ class.  The resulting cellular reciprocity is a rational $K_0$
  statement, not an object-level sheaf on the incidence-bearing cone.
\end{enumerate}
\end{theorem}

\begin{proof}
Part~\textup{(i)} is Theorems~\ref{thm:torus-Frobenius-bimodule} and
\ref{thm:finite-residual-filtration}, with the enhanced algebra model supplied
by Lemma~\ref{lem:enhanced-torus-Koszul}.  Part~\textup{(ii)} is
Lemma~\ref{lem:PGL2-functorial-lattice-condition},
Lemma~\ref{lem:PGL2-loop-orbit-quotient}, and
Theorem~\ref{thm:PGL2-residual-orbit-comparison}.  Part~\textup{(iii)} is
Lemma~\ref{lem:PGL2-full-stabilizer-support} and
Corollary~\ref{cor:PGL2-orbital-integral-realization}, including
\eqref{eq:PGL2-Weil-quotient-measure} and
\eqref{eq:PGL2-double-coset-mass}.  Part~\textup{(iv)} is
Theorem~\ref{thm:PGL2-orbital-reciprocity}.  Parts~\textup{(v)} and
\textup{(vi)} are Theorems~\ref{thm:PGL2-genuine-small-skeletonization} and
\ref{thm:PGL2-narrow-skeletonization}, respectively.
\end{proof}

\begin{definition}[Unramified absolute-$A_1$ pair]
\label{def:A1-pair}
Let $F$ be a non-Archimedean local field with finite residue field of odd
characteristic.  An \emph{unramified absolute-$A_1$ pair} $(G,K)$ consists of
a connected reductive $F$-group $G$ which is quasi-split over $F$, split over
an unramified extension, and has absolutely simple adjoint group of type
$A_1$, together with a hyperspecial subgroup
$K=\mathcal G(\cO_F)$ for a connected reductive $\cO_F$-model $\mathcal G$.
The adjective $A_1$ refers to the absolute adjoint root system.
\end{definition}

\begin{definition}[Central-neutral standard spherical family]
\label{def:standard-A1-family}
Let $p:G\to G_{\ad}$, let $\nu:G\to G_{\ab}=G/G_{\der}$, and let
$K_{\ad}=\mathcal G_{\ad}(\cO_F)$, and let
$T=G_{\gamma_m}$ with root image $T_{\rt}\subset G_{\ad}$.  A family
$\gamma_m$, $m\ge0$, is \emph{central-neutral standard spherical} if:
\begin{enumerate}[label=\textup{(\alph*)}]
\item $T_{\rt}$ is split or unramified elliptic and its adjoint quadratic
  order is $A_m=\cO_F+\varpi^m\cO_E$ for the corresponding split or
  unramified quadratic algebra $E/F$;
\item its fixed set in the reduced building is the radius-$m$ tube about the
  split apartment or the radius-$m$ ball about the elliptic fixed vertex;
\item $\nu(\gamma_m)$ is independent of $m$, belongs to $\nu(K)$, and
  \begin{equation}
   \one_K(g^{-1}\gamma_mg)
   =\one_{K_{\ad}}\bigl(p(g)^{-1}p(\gamma_m)p(g)\bigr)
   \quad(g\in G(F));
   \label{eq:central-neutral-support}
  \end{equation}
\item $G(F)$ and $T(F)$ carry the maximal-compact volume-one normalization,
  and the Weyl discriminant has valuation $2m$.
\end{enumerate}
These conditions are hypotheses, not automatic properties of an arbitrary
unramified regular family.
\end{definition}

For such a centralizer put
\begin{equation}
 \mathfrak o_{G,T}
 =\det\!\left(X_*(T/Z(G)^\circ)\otimes\Qell\right).
 \label{eq:root-orientation-line}
\end{equation}

\begin{theorem}[Theorem B: unramified absolute-$A_1$ orientation reciprocity]
\label{thm:main-A1-reciprocity}
Let $(G,K)$ be an unramified absolute-$A_1$ pair and let $\gamma_m$ be a
central-neutral standard spherical family.  Let $S/k$ be a smooth connected
parameter base over a finite field $k$, with $\ell\ne\operatorname{char}k$,
carrying the quadratic cameral cover of $T_{\rt}$, and let
$d_S$ be dimension-normalized coefficient duality.  The classification in
Section~\ref{sec:absolute-A1} is exhaustive: the shell lattice is either
saturated or type-preserving, and in the type-preserving elliptic case it has
exactly two embedding phases.  The local Haar coefficients are given by the
corresponding orbit-summed shell interpolation $I_{G,m}$.

With $J_{G,m}=Q^mI_{G,m}$ and
$W_G(z)=\sum_{m\ge0}J_{G,m}z^m$, one has
\begin{equation}
 \boxed{
 W_G(z^{-1})
 =-z\,(\mathfrak o_{G,T}\otimes d_S)W_G(z).}
 \label{eq:main-A1-reciprocity}
\end{equation}
The multiplier is independent of the central isogeny, vertex-color lattice,
and elliptic embedding phase.  Outside the saturated $PGL_2$ case the
interpolation is a rational $K_0$ class whose Frobenius traces are the actual
local orbital coefficients; no relative orbital $R_!$-push-forward or fppf
orbit-family theorem is asserted.
\end{theorem}

\begin{proof}
This is Theorems~\ref{thm:A1-local-shell-classification} and
\ref{thm:A1-orientation-reciprocity}.  The classification reduction is
Lemma~\ref{lem:A1-building-classification}, the measure comparison is
Lemma~\ref{lem:A1-shell-transfer}, the type-preserving calculation is
Proposition~\ref{prop:SL2-parity-reciprocity}, and the quotient by
$Z(G)^\circ$ is forced by Proposition~\ref{prop:U2-center-test}.
\end{proof}

For later reference, the logical status of the two theorem packages is as
follows.
\begin{description}
\item[Residual tower and torus bimodule.]
Proved at scheme level and in the selected enhanced coefficient category; no
global orbital skeleton is inferred from it.
\item[Positive-loop shell comparison.]
Proved for represented fppf orbit sheaves; this is not promoted to a
stratification of $X_m$ or $[T\backslash X_m]$.
\item[Orbital mass and reciprocity.]
Proved as a quotient-measure identity, a Frobenius-trace formula, and a
rational $K_0$ identity; no ungraded cone sheaf or convolution is asserted.
\item[Finite labelled-shell partner.]
Proved at object level on a zero-dimensional labelled skeleton; cross-shell
incidence and a Rees action remain separate problems.
\item[Height-cone shadow.]
Proved after applying the de-Tated graded character and passing to $K_0$; this
does not supply an incidence-bearing object on the height cone.
\item[Absolute-$A_1$ extension.]
Proved for Definitions~\ref{def:A1-pair} and
\ref{def:standard-A1-family} as a classification and equality of local
orbit-summed Haar masses, assembled into an explicit rational $K_0$
interpolation.  Only its saturated $PGL_2$ member has the relative
represented-shell geometry of Theorem A.
\end{description}

\subsection{Proof routes}

The arithmetic core may be read independently of most of the formal
framework.  The shortest proof route is:
\begin{enumerate}
\item the enhanced torus model and residual tower,
  Theorem~\ref{thm:torus-Frobenius-bimodule} and
  Section~\ref{sec:residual-filtration};
\item quadratic orders and represented orbit sheaves,
  Section~\ref{sec:PGL2-orbital-comparison};
\item quotient measure, trace, and orbital reciprocity,
  Corollary~\ref{cor:PGL2-orbital-integral-realization} and
  Theorem~\ref{thm:PGL2-orbital-reciprocity};
\item the finite shell Verdier partner and the separate height-character
  shadow, Section~\ref{sec:PGL2-narrow-skeletonization}.
\end{enumerate}
The short cellular section at the end is used only for the rational $K_0$
shadow; it is not a prerequisite for the orbit-shell or Haar statements.
Theorem B has a separate route: finite central lifting, classification of the
type image in the reduced building, orbit-summed measure transfer, the two
$SL_2$ elliptic phases, and the $U_2$ connected-center test.  None of those
steps is used to strengthen the relative representability assertion of
Theorem A.


\section{The enhanced torus orientation bimodule}\label{sec:torus-orientation-bimodule}
\begin{lemma}[Enhanced Kummer--Koszul algebra and its bimodule dual]
\label{lem:enhanced-torus-Koszul}
Let $p:S'\to S$ be a finite \'etale cover splitting the $d$-dimensional torus
$T$, let $\pi':T_{S'}\to S'$, and put
\[
 V_T=X^*(T_{S'})\otimes_{\mathbf Z}\Qell(-1).
\]
In the fixed enhanced six-functor formalism of the coefficient convention
there is
a functorial equivalence of augmented $E_\infty$-algebras
\begin{equation}
 K(V_T):=\operatorname{Sym}(V_T[-1])
 \xrightarrow{\ \sim\ }R\pi'_*\Qell.
 \label{eq:enhanced-torus-Koszul-algebra}
\end{equation}
Because the coefficient category is $\Qell$-linear, the underlying object is
\begin{equation}
 K(V_T)\simeq
 \bigoplus_{i=0}^d\bigwedge^iV_T[-i],
 \label{eq:enhanced-torus-exterior-object}
\end{equation}
with exterior multiplication.  Projection to the top exterior power defines
a perfect Frobenius pairing and a functorial equivalence of
$K(V_T)$-bimodules
\begin{equation}
 K(V_T)^\vee
 \xrightarrow{\ \sim\ }
 K(V_T)\otimes(\det V_T)^\vee[d].
 \label{eq:enhanced-torus-Koszul-bimodule-dual}
\end{equation}
The left and right actions in this equivalence are the two actions induced by
the $E_\infty$ multiplication, and evaluation is multiplication followed by
projection to $\det V_T[-d]$.  These maps, including evaluation and
coevaluation, are functorial for automorphisms of the character lattice.
Consequently their pullbacks to every level of the \v{C}ech nerve of $S'/S$
satisfy the descent cocycle.
\end{lemma}

\begin{proof}
For a character $\chi:T_{S'}\to\mathbb G_m$, pullback of the universal Kummer
class gives a degree-one map
$X^*(T_{S'})\otimes\Qell(-1)[-1]\to R\pi'_*\Qell$.  Additivity in $\chi$ and
the universal property of the free commutative algebra extend it to the
augmented $E_\infty$-algebra map in
\eqref{eq:enhanced-torus-Koszul-algebra}.  For a split rank-one torus, the
Kummer pro-\'etale tower computes this map by the two-term continuous Koszul
resolution, compatibly with cup products and the augmentation.  It is
therefore a quasi-isomorphism of $E_\infty$-algebras in rank one.  The
symmetric-monoidal K\"unneth equivalence identifies the rank-$d$ map with the
tensor product of these rank-one maps.  The forgetful functor from
$E_\infty$-algebras to the underlying stable category is conservative, and
the bounded constructible $t$-structure detects equivalences.  Thus this is
an object-level $E_\infty$-equivalence, not an inference made only after
passing to cohomology or $K_0$, and its underlying object is
\eqref{eq:enhanced-torus-exterior-object}.

Exterior multiplication followed by the top-wedge projection pairs
$\bigwedge^iV_T$ perfectly with $\bigwedge^{d-i}V_T$.  The resulting adjoint
map is~\eqref{eq:enhanced-torus-Koszul-bimodule-dual}; associativity and graded
commutativity give the Frobenius balance identity
\[
 \beta(xy,z)=\beta(x,yz),
 \qquad
 \beta:K(V_T)\otimes K(V_T)\longrightarrow\det V_T[-d],
\]
including the Koszul signs.  Hence the adjoint is simultaneously left- and
right-linear.  The standard
perfect-pairing evaluation--coevaluation composites give the triangle
identities.  Every construction used only characters, Kummer pullback,
multiplication, and the top exterior projection.  It is therefore functorial
for lattice isomorphisms.  On the \v{C}ech nerve the two pullbacks are related by
the descent automorphism of $X^*(T)$, and functoriality makes the triple-overlap
comparison the corresponding product.  This is exactly the descent cocycle.
\end{proof}

\begin{theorem}[Torus Frobenius bimodule]
\label{thm:torus-Frobenius-bimodule}
Let $S$ be a noetherian finite-dimensional scheme on which $\ell$ is
invertible.  Work in the fixed enhanced $\ell$-adic formalism declared in the
coefficient convention; write $\mathscr D_c(S)$ for its stable category and
$D_c^b(S,\Qell)$ for the bounded constructible homotopy category.  Let
$\pi:T\to S$ be a
smooth separated torus of constant relative dimension $d$, split by a finite
\'etale cover.  Let
\[
 \Lambda_T:=X_*(T)\otimes_{\mathbf Z}\Qell
\]
be the lisse $\Qell$-sheaf obtained from the \'etale sheaf of cocharacters,
and write
\[
 A_T^*=R\pi_*\Qell,
 \qquad
 A_T^!=R\pi_!\Qell,
 \qquad
 \mathfrak o_T=\det(\Lambda_T).
\]
Cup product is taken in the enhancement and makes $A_T^*$ an
$E_\infty$-algebra.  Put
$A_T^e=A_T^*\otimes^L(A_T^*)^{\mathrm{op}}$, and let
$(-)^\vee=R\mathcal Hom_{\Qell}(-,\Qell)$ denote coefficient-linear duality,
internal to $D_c^b(S,\Qell)$, not Verdier duality on $S$.  Write
$D(A_T^e)$ for the homotopy category of constructible $A_T^e$-modules in
$\mathscr D_c(S)$.  The projection-formula cup action makes $A_T^!$ a
dualizable object of this enhanced bimodule category.  All tensor products in
this theorem and its proof are derived, even when the
superscript $L$ is suppressed for readability.  Exterior multiplication
and relative Poincar\'e duality give a canonical equivalence there
\begin{equation}
 \boxed{
 A_T^!\simeq A_T^*\otimes\mathfrak o_T[-d].}
 \label{eq:torus-Frobenius-orientation}
\end{equation}
Equivalently,
\begin{equation}
 A_T^!\simeq (A_T^*)^\vee(-d)[-2d].
 \label{eq:torus-compact-ordinary-dual}
\end{equation}
The coefficient-duality evaluation is cup product followed by the relative
trace
\[
 A_T^*\otimes A_T^!\longrightarrow\Qell(-d)[-2d].
\]
Under~\eqref{eq:torus-Frobenius-orientation}, bimodule evaluation and
coevaluation are multiplication together with the ordinary evaluation and
coevaluation of $\mathfrak o_T[-d]$.  They satisfy the triangle identities and
exhibit $A_T^!$ as an invertible dualizing $A_T^*$-bimodule, with explicit
inverse
\begin{equation}
 (A_T^!)^{-1}=A_T^*\otimes\mathfrak o_T^\vee[d].
 \label{eq:torus-Frobenius-inverse}
\end{equation}
Here the Morita evaluation is the $A_T^*$-linear map
$A_T^!\otimes^L_{A_T^*}(A_T^!)^{-1}\to A_T^*$; it is distinct from the
coefficient pairing displayed above.  The determinant line is therefore
forced before taking $K_0$.

More explicitly, if $L=A_T^*\otimes\mathfrak o_T[-d]$ and
$L^{-1}=A_T^*\otimes\mathfrak o_T^\vee[d]$, the duality maps have domains and
codomains
\begin{align*}
 \operatorname{ev}_A&:L\otimes^L_{A_T^*}L^{-1}\longrightarrow A_T^*,\\
 \operatorname{coev}_A&:A_T^*\longrightarrow
 L^{-1}\otimes^L_{A_T^*}L.
\end{align*}
The composites
$L\to L\otimes_A L^{-1}\otimes_A L\to L$ and
$L^{-1}\to L^{-1}\otimes_A L\otimes_A L^{-1}\to L^{-1}$ are the identity.

The duality $(-)^\vee$ in this statement is not raw Verdier duality
$\D_S^{\mathrm{raw}}=R\mathcal Hom_S(-,\omega_S)$.  In particular, no
dimension-normalized coefficient involution is being used in
\eqref{eq:torus-Frobenius-orientation}; the normalized operation $d_{\Uk}$ used in
the generating series is defined separately in
\eqref{eq:normalized-U-duality} below.
\end{theorem}

\begin{proof}
The cup action
$A_T^*\otimes A_T^!\to A_T^!$ is the projection-formula form of the usual
action of ordinary cochains on compactly supported cochains.  After the
splitting cover, put
$V=X^*(T)\otimes\Qell(-1)$.  The Kummer pro-\'etale tower gives the Koszul
model for the continuous cochains of the free rank-$d$ Tate lattice precisely
as constructed in Lemma~\ref{lem:enhanced-torus-Koszul}.  The
Koszul resolution is self-dual, with determinant its top exterior generator;
therefore it gives at chain level, as a bimodule,
\[
 (A_T^*)^\vee
 \simeq A_T^*\otimes(\det V)^\vee[d]
 \simeq A_T^*\otimes\mathfrak o_T(d)[d].
\]
Here the determinant calculation, including every dual and twist, is
\[
 \det V
 =\det(X^*(T)\otimes\Qell)\otimes\Qell(-d)
 =\mathfrak o_T^\vee(-d),
 \qquad
 (\det V)^\vee=\mathfrak o_T(d).
\]
Its cohomology is the familiar exterior algebra
$\bigwedge^\bullet V$, so the chain-level determinant is compatible with the
top-wedge pairing.
Let $K(V)$ denote this finite Koszul model.  Contraction of exterior powers,
followed by projection to the top exterior generator, gives the coefficient
pairing $K(V)\otimes K(V)^\vee\to\Qell$; exterior multiplication gives the
$K(V)$-module map realizing~\eqref{eq:torus-Frobenius-orientation}.  Tensoring
that map with the ordinary evaluation
$\mathfrak o_T\otimes\mathfrak o_T^\vee\to\Qell$ gives
\[
 (A_T^*\otimes\mathfrak o_T[-d])
 \otimes^L_{A_T^*}(A_T^*\otimes\mathfrak o_T^\vee[d])
 \longrightarrow A_T^*.
\]
Insertion of the coevaluation of $\mathfrak o_T$ gives its inverse-direction
coevaluation.  The two composites are the identity because they are the
usual evaluation--coevaluation composites of an invertible line; this checks
the triangle identities in $D(A_T^e)$.  Relative Poincar\'e duality gives
\eqref{eq:torus-compact-ordinary-dual}; substitution yields
\eqref{eq:torus-Frobenius-orientation}.  Its compatibility with the left and
right $A_T^*$-actions is the projection-formula identity for cup product and
the relative trace: after pullback to the splitting cover it is exactly the
$K(V)$-balanced top-wedge pairing of
Lemma~\ref{lem:enhanced-torus-Koszul}, and the splitting pullback is
conservative.  Concretely, the relative trace pairing satisfies
\[
 \operatorname{tr}((ab)c)=\operatorname{tr}(a(bc))
\]
with the graded signs; this is the $A_T^e$-balance condition that makes the
adjoint map a bimodule morphism.  After replacing a finite splitting
cover by a finite \'etale Galois refinement, every Koszul map, multiplication,
trace, evaluation, and coevaluation above is equivariant for the splitting
group.  The \v{C}ech cocycle is the lattice-functorial cocycle proved in that
lemma.  Effective finite \'etale descent for $E_\infty$-algebras and their
module categories therefore descends the enhanced bimodule equivalence to
$S$.  Compatibility
with exterior multiplication makes it a bimodule, not merely a coefficient,
equivalence.
\end{proof}

\begin{lemma}[Compact cohomology of a torus]
\label{lem:torus-compact-cohomology}
Let $\pi:T\to S$ be a smooth $d$-dimensional torus split by a finite \'etale
cover, and put $\Lambda=X_*(T)\otimes\Qell$.  Then
\begin{equation}
 R^{2d-i}\pi_!\Qell
 \simeq \bigwedge^i\Lambda\,(i-d),
 \qquad 0\le i\le d.
 \label{eq:torus-compact-cohomology}
\end{equation}
In particular
\[
 R^d\pi_!\Qell\simeq\det\Lambda,
 \qquad
 R^{2d}\pi_!\Qell\simeq\Qell(-d).
\]
For $d=1$ its $\ell$-adic Grothendieck class is
$\Lell-[\det\Lambda]$.
\end{lemma}

\begin{proof}
After the splitting cover,
$R^i\pi_*\Qell\simeq\bigwedge^iX^*(T)\otimes\Qell(-i)$.
Smooth Poincar\'e duality identifies $R^{2d-i}\pi_!\Qell$ with the dual of
this local system twisted by $(-d)$.  The formula follows, and it descends
because all identifications are equivariant for the splitting group.
\end{proof}


\section{The universal quadratic residual filtration}
\label{sec:residual-filtration}

This section constructs the finite residual coefficient objects.  The next
section identifies them with compact positive-loop orbits in an actual
equal-characteristic $PGL_2$ affine-Springer locus.  Neither section constructs
a sheaf on the whole valued-field quotient, an infinite completed Rees object,
or an orbital convolution kernel.

The residual tower is first defined over the arithmetic parameter space
$\UZ$.  Put
\[
 \mathcal E_{\mathbf Z}=\mathcal O_{\UZ}[z]/(z^2-a),
 \qquad
 \mathcal E_{\mathbf Z}^-
 =\ker(\operatorname{Tr}_{\mathcal E_{\mathbf Z}/\mathcal O_{\UZ}}).
\]
For $r\geq1$ define the integral fppf quotient
\begin{equation}
 \mathscr Q_{r,\mathbf Z}=
 \frac{\operatorname{Res}_{\mathcal E_{\mathbf Z}[\epsilon]/(\epsilon^r)/\UZ}\mathbb G_m}
      {\operatorname{Res}_{\mathcal O_{\UZ}[\epsilon]/(\epsilon^r)/\UZ}\mathbb G_m}.
 \label{eq:Haar-residual-quotient}
\end{equation}
This is represented by a smooth affine $\UZ$-group scheme: after the finite
\'etale cover splitting $\mathcal E_{\mathbf Z}$, it is a truncated unit
group, and effective affine descent applies.

Now let $k$ be a field with $2,\ell\in k^\times$, base change to $\Uk$, and
write $U=\Uk$, $\mathcal E=\mathcal E_k$,
$\mathcal E^-=\mathcal E_{\mathbf Z}^-\otimes k$, and
$\mathscr Q_r=\mathscr Q_{r,\mathbf Z}\otimes k$.  Write
$q_r:\mathscr Q_r\to U$ and $H_r=Rq_{r,!}\Qell$.
The quotient in~\eqref{eq:Haar-residual-quotient} means the fppf sheaf
quotient; all subsequent $\ell$-adic identities live over $U=\Uk$.

\begin{theorem}[Quadratic residual tower and its orientation coefficient]
\label{thm:finite-residual-filtration}
Let $\mathbb T=\mathscr Q_1$, let
$A=R\pi_{\mathbb T,*}\Qell$, and let
$\omega_A=R\pi_{\mathbb T,!}\Qell$.
\begin{enumerate}[label=\textup{(\roman*)}]
\item The fppf quotient $\mathscr Q_1$ is canonically the norm-one torus
  $\ker(N_{\mathcal E/\mathcal O_U}:\operatorname{Res}_{\mathcal E/U}
  \mathbb G_m\to\mathbb G_m)$.
\item For $r\ge2$, reduction modulo $\epsilon^{r-1}$ is an fppf torsor
  \begin{equation}
   p_r:\mathscr Q_r\longrightarrow\mathscr Q_{r-1}
   \quad\text{under }\mathbb V(\mathcal E^-).
   \label{eq:Haar-congruence-tower}
  \end{equation}
  Consequently
  \begin{equation}
   H_r\simeq\omega_A(-(r-1))[-2(r-1)].
   \label{eq:Haar-shell-Thom}
  \end{equation}
\item Here $(m)$ denotes the ordinary Tate twist used for the chosen Weyl
  normalization; its class is $Q^m$ in the convention of
  Section~\ref{sec:quadratic-residual-series}.  After this twist and the even
  relative-dimension shift,
  \begin{equation}
   H_r(m)[2r-2]
   \simeq\omega_A(m-r+1)
   \simeq A\otimes\cK[-1](m-r+1).
   \label{eq:Haar-canonical-height-object}
  \end{equation}
  Thus the residual coefficient height is $s=m-r+1$.
\item The motivic and $\ell$-adic class identities are, respectively,
  \begin{equation}
   [\mathscr Q_{r,\mathbf Z}]
   =(\Lmot-\kappamot)\Lmot^{r-1}
   \quad\text{in }K_0(\operatorname{Var}_{\UZ}),
   \label{eq:residual-layer-motivic-class}
  \end{equation}
  and
  \begin{equation}
   [Rq_{r,!}\Qell]
   =(\Lell-\kappaell)\Lell^{r-1}
   \quad\text{in }K_0(D_c^b(\Uk,\Qell)).
   \label{eq:residual-layer-class}
  \end{equation}
\end{enumerate}
\end{theorem}

\begin{proof}
Because $2$ is invertible, the quadratic algebra is \'etale and its
trace-zero summand is a line bundle.  The morphism
$x\mapsto x/\sigma(x)$ identifies the first quotient with the norm-one torus
after the splitting cover; the identification descends and proves (i).
For $r\ge2$, the kernel of reduction is
\[
 \frac{1+\epsilon^{r-1}\mathcal E}
      {1+\epsilon^{r-1}\mathcal O_U}
 \simeq \mathcal E/\mathcal O_U
 \simeq\mathcal E^-.
\]
This calculation is equivariant for the quadratic descent action, so the
reduction morphism is the asserted vector-line torsor, proving
\eqref{eq:Haar-congruence-tower}.  Compact homotopy invariance gives
$Rp_{r,!}\Qell\simeq\Qell(-1)[-2]$.  Iteration proves
\eqref{eq:Haar-shell-Thom}.

The quadratic deck involution acts by $-1$ on $X_*(\mathbb T)$, so
$\det(X_*(\mathbb T)\otimes\Qell)$ is precisely the sign local system $\cK$.
Theorem~\ref{thm:torus-Frobenius-bimodule} in dimension one therefore gives
$\omega_A\simeq A\otimes\cK[-1]$, proving
\eqref{eq:Haar-canonical-height-object}.  Finally the vector-torsor tower
over $\UZ$ gives~\eqref{eq:residual-layer-motivic-class}, using
$[\mathbb T_{\mathbf Z}]=\Lmot-\kappamot$; the latter follows by subtracting
the quadratic boundary $\widetilde\UZ$ from the split conic compactification.
Compactly supported realization gives~\eqref{eq:residual-layer-class}, and
Lemma~\ref{lem:torus-compact-cohomology} identifies the realized first layer
as $\Lell-\kappaell$.
\end{proof}

\begin{remark}[The residual filtration before orbital comparison]
For a fixed $m$, the disjoint union
\[
 \mathscr B_m=U\amalg\coprod_{1\le r\le m}\mathscr Q_r,
 \qquad
 F_{\le r}\mathscr B_m=U\amalg\coprod_{1\le j\le r}\mathscr Q_j
\]
is a finite split filtration by open-and-closed pieces, with positive
coefficient layers $H_r$.  Theorem~\ref{thm:PGL2-residual-orbit-comparison}
will identify each piece with an actual compact orbit shell.  The split union
itself still creates no geometric transition maps between different orbital
depths.  In particular
the reduction morphisms $\mathscr Q_r\to\mathscr Q_{r-1}$ are not a Rees
action, and no inverse system or completion in $m$ is asserted.
\end{remark}

\begin{proposition}[The residual Kummer coefficient]
\label{prop:residual-kummer-bimodule}
The actual compact coefficient complex of the first residual layer is
$\omega_A$.  With $A^e=A\otimes^LA^{\mathrm{op}}$ and $D(A^e)$ interpreted
in the enhanced sense fixed in
Theorem~\ref{thm:torus-Frobenius-bimodule}, one has
\[
 \omega_A\simeq A\otimes\cK[-1].
\]
This is an object-level Frobenius-bimodule identity for the residual torus.
After Theorem~\ref{thm:PGL2-residual-orbit-comparison} it is also the vertical
coefficient identity for the first actual compact orbit shell.  By itself it
is not a skeletonization theorem or a completed Rees self-duality.
\end{proposition}


\section{From residual quotients to actual
\texorpdfstring{$PGL_2$}{PGL2} compact orbital shells}
\label{sec:PGL2-orbital-comparison}

We now work in equal characteristic.  The purpose of this section is to close
one precise geometric gap: the schemes $\mathscr Q_r$ will be identified with
positive-loop torus orbits inside an actual group affine-Springer locus.  The
affine Grassmannian is used in its standard fpqc-sheaf sense; its
representability and basic orbit geometry are recalled in
\cite{PappasRapoport,ZhuAffineGrassmannians}.  We do not identify the whole
quotient stack of the affine-Springer locus with a finite-type scheme.

Put $\mathcal O_R=R[[\varpi]]$ and $F_R=R((\varpi))$ for a $U$-algebra $R$.
In the basis $(1,z)$ of $\mathcal E$, use the regular representation
\begin{equation}
 \iota(z)=
 \begin{pmatrix}0&a\\1&0\end{pmatrix},
 \qquad \iota(z)^2=a.
 \label{eq:PGL2-quadratic-embedding}
\end{equation}
Let $G=PGL(\mathcal E)\simeq PGL_{2,U}$ and let
\begin{equation}
 \gamma_m=\left[1+\varpi^m\iota(z)\right]
 =\left[\begin{pmatrix}1&a\varpi^m\\\varpi^m&1\end{pmatrix}\right]
 \in G(\mathcal O_U((\varpi))),\qquad m\ge0.
 \label{eq:PGL2-universal-gamma}
\end{equation}
Its determinant is $1-a\varpi^{2m}$.  This is a unit for $m>0$, and for
$m=0$ the same statement is exactly why $(1-a)^{-1}$ was included in $U$.
By~\eqref{eq:companion-integral-conjugacy}, this is the same regular
semisimple conjugacy family as the companion presentation, now written in the
integral quadratic basis used throughout the orbit-sheaf proof.

The centralizer of $\gamma_m$ is the unramified torus
\begin{equation}
 \mathbb T=(\operatorname{Res}_{\mathcal E/U}\mathbb G_m)/\mathbb G_m.
 \label{eq:PGL2-centralizer-torus}
\end{equation}
Write $\operatorname{Gr}_G=LG/L^+G$ as an fpqc quotient sheaf.  Intrinsically,
an $R$-point is a $G$-torsor $\mathcal P$ on
$D_R=\Spec R[[\varpi]]$ together with a trivialization on
$D_R^\times=\Spec R((\varpi))$.  The trivialization lets $\gamma_m$ act on
$\mathcal P|_{D_R^\times}$.  Define the group affine-Springer subfunctor by
\begin{equation}
 X_m(R)=\{(\mathcal P,\alpha):
 \alpha^{-1}\gamma_m\alpha\text{ extends to an automorphism of }\mathcal P
 \text{ on }D_R\}.
 \label{eq:PGL2-group-affine-Springer}
\end{equation}
When a section is represented by $g\in LG(R)$ this is the familiar condition
$g^{-1}\gamma_mg\in L^+G(R)$.  The torsor formulation is required over a
general base: we do not assume that a $PGL_2$ loop admits a $GL_2$ lift after
a cover pulled back from $\Spec R$.  Affine-Springer fibers and their relation to orbital
integrals provide the broader context in~\cite{GKM}.

\begin{lemma}[Projective-lattice descent and the liftable locus]
\label{lem:PGL2-functorial-lattice-condition}
Let $R$ be a $U$-algebra.
\begin{enumerate}[label=\textup{(\roman*)}]
\item The condition in~\eqref{eq:PGL2-group-affine-Springer} is invariant under
  isomorphism of the pair $(\mathcal P,\alpha)$ and is compatible with fpqc
  base change.  Hence it defines a sub-fpqc-sheaf of
  $\operatorname{Gr}_G$ without choosing a $GL_2$ lift.
\item Suppose $(\mathcal P,\alpha)$ admits a lift to a rank-two projective
  $R[[\varpi]]$-module $L$ with a trivialization over $R((\varpi))$.
  Then membership in $X_m(R)$ is equivalent to
  \[
   \widetilde\gamma_mL=L,
  \]
  and also to stability under the order
  \[
   \mathcal A_{m,R}=R[[\varpi]]+\varpi^m
   (\mathcal E\otimes_{\mathcal O_U}R)[[\varpi]]
   =R[[\varpi]][\widetilde\gamma_m].
  \]
  Tensoring $L$ by a line bundle, or changing a local matrix lift by a
  scalar, does not change either condition.
\item Every positive-loop orbit section $u[\mathcal A_r]$ used below has the
  explicit lift $u\mathcal A_{r,R}$.  Every field-valued point of
  $\operatorname{Gr}_G$ is liftable as well.  Thus the arbitrary-base orbit
  theorem and the field-valued quadratic-order classification both lie in
  the liftable locus, although no lifting assertion is made for an arbitrary
  $R$-point of the whole affine Grassmannian.
\end{enumerate}
\end{lemma}

\begin{proof}
The Beauville--Laszlo torsor description of the affine Grassmannian is fpqc
functorial.  Extension of a generic automorphism across $D_R$ is invariant
under descent isomorphisms, proving~\textup{(i)}.

For~\textup{(ii)}, work fpqc locally on $D_R$ where the projective module is
free and choose a matrix representative $g$.  Put
$h=g^{-1}\widetilde\gamma_mg$.  If the class of $h$ is integral in $PGL_2$,
there is locally a scalar $c$ with $ch\in GL_2(R[[\varpi]])$.  Since
$\operatorname{tr}(h)=2$ and $2\in R^\times$, one has
$c=\operatorname{tr}(ch)/2\in R[[\varpi]]$.  Moreover
$\det(h)=1-a\varpi^{2m}$ is a unit, so
$\det(ch)=c^2\det(h)$ being a unit forces
$c\in R[[\varpi]]^\times$.  Hence $h$ itself is integral, which is exactly
$\widetilde\gamma_mL=L$.  Now
$\widetilde\gamma_m-1=\varpi^m\iota(z)$, and
\[
 \widetilde\gamma_m^{-1}
 =(1-a\varpi^{2m})^{-1}(1-\varpi^m\iota(z)).
\]
Therefore preservation by $\widetilde\gamma_m$ is equivalent to
$\mathcal A_{m,R}$-stability.  The construction descends because it was
phrased in terms of the projective module and its endomorphism.

For~\textup{(iii)}, multiplication by $u$ gives the displayed module lift.
For a field-valued affine-Grassmannian point, the corresponding $PGL_2$-torsor
on $\mathcal O_F$ is trivial over $F$.  Its Brauer class is therefore in the
kernel of $\operatorname{Br}(\mathcal O_F)\to\operatorname{Br}(F)$, which is
zero for the regular discrete valuation ring $\mathcal O_F$.  The remaining
vector-bundle lift is free because $\mathcal O_F$ is local.  After choosing
that trivialization, Hilbert~90 applied to
$GL_2(F)\to PGL_2(F)\to H^1(F,\mathbb G_m)$ gives a matrix representative.
This proves the two stated lifting cases.
\end{proof}

For $r\ge0$ put
\begin{equation}
 \mathcal A_{r,R}=R[[\varpi]]+\varpi^r
        (\mathcal E\otimes_{\mathcal O_U}R)[[\varpi]].
 \label{eq:PGL2-order-lattice}
\end{equation}
This is a rank-two $R[[\varpi]]$-lattice and a quadratic order.  Let
$L^+\mathbb T$ denote the fppf positive-loop sheaf
\begin{equation}
 L^+\mathbb T(R)=
 \frac{(\mathcal E\otimes R)[[\varpi]]^\times}
      {R[[\varpi]]^\times},
 \label{eq:PGL2-positive-loop-torus}
\end{equation}
and let $L^+\mathbb T_r$ be the stabilizer sheaf of
$[\mathcal A_r]$.

\begin{remark}[Meaning of ``compact shell'']
The adjective \emph{compact} refers to the orbit of the maximal compact
subgroup $T(\mathcal O)$, or equivalently to the positive-loop orbit.  It does
not mean that the representing algebraic $U$-scheme is proper: after the
quadratic splitting cover, $\mathscr Q_r$ contains a $\mathbb G_m$ factor.
\end{remark}

\begin{lemma}[Local quadratic-order classification]
\label{lem:PGL2-local-order-classification}
Let $F$ be a complete discretely valued field of residue characteristic
different from two, and let $E/F$ be a split or unramified quadratic
\'etale algebra with maximal order $\mathcal O_E$.  Put
$A_m=\mathcal O_F+\varpi^m\mathcal O_E$.  Every $A_m$-stable full
$\mathcal O_F$-lattice in $E$, modulo multiplication by $E^\times$, has a
unique representative
\[
 A_r=\mathcal O_F+\varpi^r\mathcal O_E,
 \qquad 0\le r\le m.
\]
Its multiplier ring is $A_r$.
\end{lemma}

\begin{proof}
Every fractional $\mathcal O_E$-ideal is principal.  Scale a lattice $L$ so
that $\mathcal O_EL=\mathcal O_E$ and hence $L\subset\mathcal O_E$.  In the
field case $L$ contains a unit.  In the split case the two projections of
$L/\varpi L$ are nonzero; a vector space over a field cannot be the union of
its two coordinate-axis intersections, so $L$ again contains an element that
is a unit in both factors.  Scaling by that unit puts $1$ in $L$.

Choose $\omega$ with
$\mathcal O_E=\mathcal O_F\oplus\mathcal O_F\omega$.  The image of $L$ in
$\mathcal O_E/\mathcal O_F$ is
$\varpi^r\mathcal O_F\bar\omega$ for a unique $r\ge0$.  Subtracting a scalar
from a lift of its generator gives
$L=\mathcal O_F+\varpi^r\mathcal O_E=A_r$.  Stability under $A_m$ is
equivalent to $\varpi^m\mathcal O_E\subset A_r$, hence to $r\le m$.
Finally $xA_r\subset A_r$ implies $x=x\cdot1\in A_r$, and the converse holds
because $A_r$ is a ring.  Thus the multiplier ring is $A_r$ and determines
$r$.
\end{proof}

\begin{lemma}[The positive-loop orbit is the residual quotient]
\label{lem:PGL2-loop-orbit-quotient}
For $r\ge1$ there are canonical isomorphisms of fppf sheaves over $U$
\begin{align}
 L^+\mathbb T_r
 &\simeq
 \mathcal A_r^\times/\mathcal O_U[[\varpi]]^\times,
 \label{eq:PGL2-loop-stabilizer}\\
 L^+\mathbb T/L^+\mathbb T_r
 &\simeq
 \frac{\mathcal E[[\varpi]]^\times}
 {\mathcal O_U[[\varpi]]^\times(1+\varpi^r\mathcal E[[\varpi]])}
 \simeq \mathscr Q_r.
 \label{eq:PGL2-loop-quotient-Qr}
\end{align}
For $r=0$ the quotient is the point $U=\mathscr Q_0$.
For every $m\ge r$, the induced orbit morphism
\[
 L^+\mathbb T/L^+\mathbb T_r\longrightarrow X_m
\]
is a monomorphism of fppf sheaves.
\end{lemma}

\begin{proof}
For a $U$-algebra $R$, an element of $L^+\mathbb T(R)$ is fppf locally
represented by $u\in(\mathcal E\otimes R)[[\varpi]]^\times$.  It fixes the
homothety class of $\mathcal A_{r,R}$ precisely when there is
$c\in R((\varpi))^\times$ with
$u\mathcal A_{r,R}=c\mathcal A_{r,R}$.  Set $a=c^{-1}u$.  The displayed equality says
$a\mathcal A_{r,R}=\mathcal A_{r,R}$; since $1\in\mathcal A_{r,R}$, applying
the same argument to $a^{-1}$ gives
$a\in\mathcal A_{r,R}^\times$.  Hence
$c=ua^{-1}$ belongs to
$(\mathcal E\otimes R)[[\varpi]]^\times$.  Because
$\mathcal E\otimes R$ is finite faithfully flat over $R$, scalar Laurent
series satisfy
\[
 R((\varpi))^\times\cap
 (\mathcal E\otimes R)[[\varpi]]^\times
 =R[[\varpi]]^\times;
\]
this is the equalizer form of faithfully flat descent, applied also to the
inverse.  Thus $c\in R[[\varpi]]^\times$ fppf locally, and the homothety may
be absorbed by the scalar subgroup.  Conversely every unit of the order
preserves it.  This proves~\eqref{eq:PGL2-loop-stabilizer} as an equality of
fppf stabilizer sheaves, including the scalar quotient.
For every $U$-algebra $R$,
\[
 \mathcal A_{r,R}^\times
 =R[[\varpi]]^\times
  (1+\varpi^r(\mathcal E\otimes R)[[\varpi]]).
\]
Indeed the reduction of a unit of $\mathcal A_{r,R}$ modulo
$\varpi^r\mathcal E[[\varpi]]$ is a scalar unit, and division by a lift of
that scalar leaves an element of the displayed congruence subgroup.
Reduction modulo $\varpi^r$ therefore identifies the middle quotient in
\eqref{eq:PGL2-loop-quotient-Qr} with
\[
 \frac{((\mathcal E\otimes R)[\varpi]/(\varpi^r))^\times}
      {(R[\varpi]/(\varpi^r))^\times},
\]
which is exactly~\eqref{eq:Haar-residual-quotient} after
$\epsilon\mapsto\varpi$.  Every construction commutes with fppf base change,
so the calculation for arbitrary $R$ is an isomorphism after fppf
sheafification.  Finally, if two sections $u_1,u_2$ have the same image in
$X_m$, then fppf locally
$u_2^{-1}u_1$ fixes $[\mathcal A_r]$ and hence belongs to
$L^+\mathbb T_r$ by the stabilizer calculation.  They therefore define the
same quotient section.  This proves that the orbit map is a monomorphism of
fppf sheaves, not only an injection on field-valued points.
\end{proof}

\begin{theorem}[Represented positive-loop orbit-shell comparison for the universal $PGL_2$ family]
\label{thm:PGL2-residual-orbit-comparison}
For $m\ge0$ and $0\le r\le m$, the order lattice $\mathcal A_r$ is fixed by
$\gamma_m$.  The orbit morphism
\begin{equation}
 \theta_{m,r}:\mathscr Q_r
 \xrightarrow{\ \sim\ }
 L^+\mathbb T\cdot[\mathcal A_r]
 \longrightarrow X_m,
 \qquad [u]\longmapsto[u\mathcal A_r],
 \label{eq:PGL2-orbit-morphism}
\end{equation}
identifies $\mathscr Q_r$ with the fppf positive-loop orbit sheaf.  In
particular that orbit sheaf is represented by a smooth finite-type $U$-scheme,
and the induced morphism of fpqc sheaves to
$X_m\subset\operatorname{Gr}_G$ is a monomorphism.  No representability
statement for all of $X_m$ is used here.

Write $\mathscr O_{m,r}$ for this represented orbit and set
\begin{equation}
 \mathscr S_m^{\mathrm{orb}}
 =\coprod_{r=0}^m\mathscr O_{m,r}.
 \label{eq:PGL2-actual-shell-union}
\end{equation}
By definition of this disjoint union, the componentwise orbit isomorphisms
assemble into
\begin{equation}
 \Theta_m:
 U\amalg\coprod_{r=1}^m\mathscr Q_r
 \xrightarrow{\ \sim\ }\mathscr S_m^{\mathrm{orb}}.
 \label{eq:PGL2-residual-orbital-isomorphism}
\end{equation}
This displayed map is a convenient bookkeeping isomorphism, not an additional
decomposition theorem for $X_m$ or $[T\backslash X_m]$.
For every field-valued point $x\to U$, with
$F_x=k(x)((\varpi))$ and $E_x/F_x$ the induced split or quadratic unramified
algebra, the double quotient
\[
 T_x(F_x)\backslash X_{m,x}(k(x))
\]
has exactly the representatives $[A_r]$, $0\le r\le m$.  Thus
$\mathscr O_{m,r}$ is a compact fundamental shell for the corresponding
orbital double coset.  The theorem does not assert that
$[T\backslash X_m]$ is the disjoint union of the $\mathscr Q_r$.
\end{theorem}

\begin{proof}
Since $\iota(z)\mathcal E=\mathcal E$,
\[
 \varpi^m\iota(z)\mathcal A_r\subset\mathcal A_r
 \quad\Longleftrightarrow\quad m\ge r.
\]
The inverse of $1+\varpi^m\iota(z)$ is
$(1-a\varpi^{2m})^{-1}(1-\varpi^m\iota(z))$, so containment is equality.
Hence $[\mathcal A_r]\in X_m$ and the centralizer action preserves $X_m$.
The stabilizer calculation in Lemma~\ref{lem:PGL2-loop-orbit-quotient} proves
that the fppf orbit is $\mathscr Q_r$ and that its orbit map is a monomorphism.
Theorem~\ref{thm:finite-residual-filtration} supplies smooth finite-type
representability.  The componentwise isomorphisms assemble tautologically to
\eqref{eq:PGL2-residual-orbital-isomorphism}; multiplier rings distinguish the
field-valued double-coset representatives indexed by different $r$.

Over a field-valued point, a coset in $X_m$ is a homothety class of
$A_m=\mathcal O_{F_x}[\widetilde\gamma_m]$-stable lattices.  Lemma
\ref{lem:PGL2-local-order-classification} gives precisely the representatives
$A_r$, and their stabilizers are
$A_r^\times/\mathcal O_{F_x}^\times$.  In the anisotropic case
$T_x(F_x)=T_x(\mathcal O_{F_x})$; in the split case translations along the
apartment are supplied by $T_x(F_x)/T_x(\mathcal O_{F_x})$.  In both cases
the compact $L^+T$-orbit is a fundamental shell whose cardinality represents
the inverse stabilizer volume of the double coset.  This order calculation is
the equal-characteristic form of the rank-one tube/ball decomposition familiar
from the explicit orbital formulas in~\cite[Sections~5.8--5.9]{KottwitzOrbital}.
\end{proof}

\begin{lemma}[Full torus stabilizer and orbital support]
\label{lem:PGL2-full-stabilizer-support}
Let $x\to U$ be field-valued, put $F=F_x$, $K=PGL_2(\mathcal O_F)$, and
choose $g_r$ with $g_r\mathcal O_F^2=A_{r,x}$.  Then
\begin{equation}
 \operatorname{Stab}_{T_x(F)}([A_{r,x}])
 =T_x(F)\cap g_rKg_r^{-1}
 =A_{r,x}^\times/\mathcal O_F^\times.
 \label{eq:PGL2-full-torus-stabilizer}
\end{equation}
Moreover the support of the spherical orbital integrand is the disjoint union
\begin{equation}
 \{g\in PGL_2(F):g^{-1}\gamma_{m,x}g\in K\}
 =\coprod_{r=0}^m T_x(F)g_rK.
 \label{eq:PGL2-orbital-support-decomposition}
\end{equation}
\end{lemma}

\begin{proof}
Write $T_x(F)=E_x^\times/F^\times$.  If the class of $e\in E_x^\times$
stabilizes $[A_{r,x}]$, then
$eA_{r,x}=cA_{r,x}$ for some $c\in F^\times$.  Hence
$c^{-1}e$ is a unit of the multiplier ring of $A_{r,x}$, which is
$A_{r,x}$ by Lemma~\ref{lem:PGL2-local-order-classification}.  Conversely
every such unit stabilizes the lattice.  The kernel of
$A_{r,x}^\times\to E_x^\times/F^\times$ is
$A_{r,x}^\times\cap F^\times=\mathcal O_F^\times$, proving
\eqref{eq:PGL2-full-torus-stabilizer}.  Finally the set on the left of
\eqref{eq:PGL2-orbital-support-decomposition}, modulo right translation by
$K$, is $X_{m,x}$.  The field-valued double-quotient classification in
Theorem~\ref{thm:PGL2-residual-orbit-comparison} gives exactly the classes
$[A_{r,x}]$; lifting those classes by the chosen $g_r$ proves the disjoint
support decomposition.
\end{proof}

\begin{remark}[Orbit shells are not closure strata of the whole fiber]
The schemes $\mathscr O_{m,r}$ are compact fundamental representatives for
orbital double cosets, not a claimed stratification of all of $X_m$.  This
distinction is already visible after the quadratic splitting cover: the
closure of the $\mathbb G_m$ direction in an ambient Schubert variety acquires
boundary points translated along the split apartment, whereas
$\mathscr S_m^{\mathrm{orb}}$ retains one compact fundamental shell per
$T(F)$-double coset.  Consequently neither closure relations nor transition
maps may be read off from the disjoint union
\eqref{eq:PGL2-actual-shell-union}.
\end{remark}

\begin{corollary}[Geometric realization of the spherical orbital integral]
\label{cor:PGL2-orbital-integral-realization}
Let $x\in U$ be a closed point with finite residue field $k_x$, put
$F_x=k_x((\varpi))$, and normalize Haar measures by
\[
 \operatorname{vol}(PGL_2(k_x[[\varpi]]))=1,
 \qquad \operatorname{vol}(T_x(k_x[[\varpi]]))=1.
\]
The quotient measure $d\dot g$ is the one characterized by Weil integration:
for every compactly supported locally constant $f$,
\begin{equation}
 \int_{PGL_2(F_x)}f(g)\,dg
 =\int_{T_x(F_x)\backslash PGL_2(F_x)}
   \int_{T_x(F_x)}f(tg)\,dt\,d\dot g.
 \label{eq:PGL2-Weil-quotient-measure}
\end{equation}
Let $\operatorname{Frob}_x$ denote geometric Frobenius.  In the last line
below, $\Tr$ denotes the derived supertrace
$\sum_i(-1)^i\operatorname{Tr}(\operatorname{Frob}_x\mid\mathcal H^i)$.
Then
\begin{align}
 \operatorname{OI}_m(x)
 &:={\int}_{T_x(F_x)\backslash PGL_2(F_x)}
 \one_{PGL_2(k_x[[\varpi]])}(g^{-1}\gamma_{m,x}g)\,d\dot g
 \notag\\
 &=\sum_{r=0}^m
   [T_x(k_x[[\varpi]]):A_{r,x}^\times/k_x[[\varpi]]^\times]
 =\#\mathscr S_{m,x}^{\mathrm{orb}}(k_x)
 \notag\\
 &=\Tr\!\left(\operatorname{Frob}_x\mid
   (Rq^{\mathrm{orb}}_{m,!}\Qell)_{\bar x}\right),
 \label{eq:PGL2-orbital-trace}
\end{align}
where $q_m^{\mathrm{orb}}:\mathscr S_m^{\mathrm{orb}}\to U$ is the structure
map.
\end{corollary}

\begin{proof}
Put $K_x=PGL_2(k_x[[\varpi]])$ and choose $g_r$ with
$g_rk_x[[\varpi]]^2=A_{r,x}$.  By
\eqref{eq:PGL2-orbital-support-decomposition}, the support of the
characteristic function is precisely the disjoint union
$\coprod_{r=0}^mT_x(F_x)g_rK_x$.  Formula~\eqref{eq:PGL2-Weil-quotient-measure}
therefore gives
\begin{equation}
 \operatorname{vol}\bigl(T_x(F_x)\backslash
 T_x(F_x)g_rK_x\bigr)
 =\frac{\operatorname{vol}(K_x)}
 {\operatorname{vol}(T_x(F_x)\cap g_rK_xg_r^{-1})}.
 \label{eq:PGL2-double-coset-mass}
\end{equation}
The full-torus stabilizer calculation
\eqref{eq:PGL2-full-torus-stabilizer}, not merely the positive-loop
stabilizer calculation, identifies the denominator with
$A_{r,x}^\times/\mathcal O_{F_x}^\times$, as a subgroup of
\[
 T_x(\mathcal O_{F_x})
 =\mathcal O_{E_x}^\times/\mathcal O_{F_x}^\times.
\]
Since the latter has volume one,~\eqref{eq:PGL2-double-coset-mass} is exactly
\[
 [T_x(\mathcal O_{F_x}):
   A_{r,x}^\times/\mathcal O_{F_x}^\times].
\]
In the unramified-field case $T_x(F_x)=T_x(\mathcal O_{F_x})$.  In the split
case $T_x(F_x)/T_x(\mathcal O_{F_x})$ is the apartment-translation lattice;
it is already divided out on the left of the quotient, while the stabilizer
in~\eqref{eq:PGL2-double-coset-mass} remains the same compact subgroup.
Thus no additional split-torus volume factor occurs.

Finally, the fiber of $\mathscr Q_r$ is the fppf quotient of the unit groups
of the two finite $k_x$-algebras
$E_x[\epsilon]/(\epsilon^r)$ and
$k_x[\epsilon]/(\epsilon^r)$.  The obstruction to replacing quotient-sheaf
$k_x$-points by the naive unit quotient lies in
$H^1(k_x,(k_x[\epsilon]/(\epsilon^r))^\times)$, which is the Picard group of
this finite local Artin algebra and is zero.  Equivalently this is the
truncated-unit form of Hilbert 90.  Hence
$\mathscr Q_r(k_x)$ is the ordinary unit quotient and its cardinality is the
same index.  Sum over the representatives from
Theorem~\ref{thm:PGL2-residual-orbit-comparison} and apply the
Grothendieck--Lefschetz trace formula.
\end{proof}

\begin{proposition}[Independent Bruhat--Tits verification of the shell support]
\label{prop:PGL2-building-shell-verification}
At a field-valued point $x\to U$, let $\mathcal B_x$ be the reduced
Bruhat--Tits tree of $PGL_2(F_x)$.  If $T_x$ is split, let $\mathcal A_x$ be
its apartment; if $T_x$ is unramified anisotropic, let $v_x$ be its fixed
vertex.  Then the fixed set of $\gamma_{m,x}$ is, respectively, the closed
radius-$m$ tube around $\mathcal A_x$ or the closed radius-$m$ ball around
$v_x$.  Modulo $T_x(F_x)$ it has one shell at each distance
$0\le r\le m$, and the stabilizer of the distance-$r$ shell is
$A_{r,x}^{\times}/\mathcal O_{F_x}^{\times}$.

Consequently the building calculation independently recovers both the
support decomposition~\eqref{eq:PGL2-orbital-support-decomposition} and the
termwise masses in~\eqref{eq:PGL2-orbital-trace}.
\end{proposition}

\begin{proof}
After the unramified splitting extension, the two eigenvalues of
$\gamma_{m,x}$ are $1\pm\varpi^m\sqrt a$.  Their difference has valuation
$m$.  In a lattice chart transverse to the eigenline apartment, moving one
edge away divides one off-diagonal integrality condition by $\varpi$; hence
exactly the vertices of transverse distance at most $m$ are fixed.  Galois
descent gives the tube in the split case and the ball about the unique
anisotropic fixed vertex in the nonsplit case.  Apartment translations are
already supplied by $T_x(F_x)$, so there is one orbit for each distance.

The lattice at distance $r$ may be represented by
$A_{r,x}=\mathcal O_{F_x}+\varpi^r\mathcal O_{E_x}$.  Its multiplier ring is
$A_{r,x}$, and the same calculation as in
Lemma~\ref{lem:PGL2-full-stabilizer-support} identifies its torus stabilizer
with $A_{r,x}^{\times}/\mathcal O_{F_x}^{\times}$.  Weil's quotient formula
then gives the stated masses.  This proof uses the metric geometry of the
tree rather than the fppf orbit representation, so it is a consistency check
independent of the construction of $\mathscr Q_r$.
\end{proof}

\begin{theorem}[Universal $PGL_2$ compact-orbital reciprocity]
\label{thm:PGL2-orbital-reciprocity}
Let
$\mathcal I_m^{\mathrm{orb}}=Rq^{\mathrm{orb}}_{m,!}\Qell$ and let
$I_m^{\mathrm{orb}}=[\mathcal I_m^{\mathrm{orb}}]$ be the raw compact-orbital
coefficient.  Define its Weyl-normalized coefficient by
$J_m^{\mathrm{orb}}=Q^mI_m^{\mathrm{orb}}$.  The group Weyl
discriminant of~\eqref{eq:PGL2-universal-gamma} is as follows; in the absolute
value statement $x$ is a finite-field point and $q_x=|k_x|$:
\begin{equation}
 D^G(\gamma_m)
 =\frac{-4a\varpi^{2m}}{1-a\varpi^{2m}},
 \qquad |D^G(\gamma_{m,x})|_x^{1/2}=q_x^{-m},
 \quad q_x=|k_x|,
 \label{eq:PGL2-Weyl-discriminant}
\end{equation}
so the motivic Weyl factor is $\Lmot^{-m}$ and its $\ell$-adic realization is
$Q^m$.  In the localized
Tate--Kummer Grothendieck ring,
\begin{align}
 J_m^{\mathrm{orb}}
 &=1+(1-\kappaell)\sum_{s=1}^mQ^s,
 \label{eq:PGL2-normalized-coefficient}\\
 W_{PGL_2}(z)
 :=\sum_{m\ge0}J_m^{\mathrm{orb}}z^m
 &=\frac{1-\kappaell Qz}{(1-z)(1-Qz)},
 \label{eq:PGL2-orbital-series}
\end{align}
Moreover,
\begin{equation}
 \boxed{W_{PGL_2}(z^{-1})
 =-\kappaell z\,d_{\Uk}(W_{PGL_2}(z)),\qquad
 \kappaell=[\det(X_*(\mathbb T)\otimes\Qell)]=[\cK].}
 \label{eq:PGL2-actual-orbital-reciprocity}
\end{equation}
At every finite-field point, the geometric-Frobenius derived supertrace of the coefficient in
\eqref{eq:PGL2-orbital-series} is the Weyl-normalized actual spherical orbital
integral in~\eqref{eq:PGL2-orbital-trace}.
\end{theorem}

\begin{proof}
The two eigenvalues of $\gamma_m$ are
$1\pm\varpi^m\sqrt a$, which gives~\eqref{eq:PGL2-Weyl-discriminant}; its
denominator is a unit on $U$.  This normalization is also the intrinsic
root-space Jacobian, not a factor chosen to make the series reciprocal.
Indeed, in the standard basis $H,E,F$ of
$\mathfrak{pgl}_2\simeq\mathfrak{sl}_2$, the torus direction is
$X_a=aE+F$ and a transverse direction is $Y_a=aE-F$; one has
$\det(X_a,H,Y_a)=2a\in\mathcal O_U^\times$.  Thus the local slice contributes
no additional valuation, while the determinant of the orbit differential on
the two root spaces is $D^G(\gamma_m)$.

The componentwise represented-orbit
isomorphisms assembled in
\eqref{eq:PGL2-residual-orbital-isomorphism} identify
$\mathcal I_m^{\mathrm{orb}}$ with the direct sum of the residual complexes.
Theorem~\ref{thm:finite-residual-filtration} and the calculation in
Section~\ref{sec:quadratic-residual-series} give
\eqref{eq:PGL2-normalized-coefficient}--\eqref{eq:PGL2-actual-orbital-reciprocity}.
Corollary~\ref{cor:PGL2-orbital-integral-realization} gives the final trace
statement.
\end{proof}


\section{The \texorpdfstring{$PGL_2$}{PGL2} compact-shell generating class}
\label{sec:quadratic-residual-series}

Let $U=\Uk=\Spec k[a,a^{-1},(1-a)^{-1}]$ be the realization base of the quadratic tower and
of the $PGL_2$ family in Sections~\ref{sec:residual-filtration} and
\ref{sec:PGL2-orbital-comparison}, and put
\[
 K_{\Uk}=K_0(D_c^b(\Uk,\Qell))[\Lell^{-1}],
 \qquad
 \Lell=[\Qell(-1)[-2]],\qquad Q=\Lell^{-1}.
\]
For a complex $M$ write $[M]$ for its class; for a finite-type map
$f:Y\to U$ abbreviate $[Y]=[Rf_!\Qell]$.  Let
$\widetilde U=\Spec k[b,b^{-1},(1-b^2)^{-1}]\to U$, $b^2=a$, be the
quadratic cover.
Define $\cK$ by
$R\widetilde p_*\Qell\simeq\Qell\oplus\cK$, with the deck involution acting
by $-1$ on $\cK$, and put $\kappaell=[\cK]$.  Raw Verdier duality on the smooth
curve $U$ is
\[
 \D_U^{\mathrm{raw}}(M)
 =R\mathcal Hom_U(M,\Qell(1)[2]).
\]
Define the dimension-normalized coefficient duality by
\begin{equation}
 \mathsf d_{\Uk}(M)
 :=\D_U^{\mathrm{raw}}(M)\otimes^L\Qell(-1)[-2],
 \qquad d_{\Uk}([M]):=[\mathsf d_{\Uk}(M)].
 \label{eq:normalized-U-duality}
\end{equation}
Thus $\mathsf d_{\Uk}$ fixes the constant coefficient, unlike raw Verdier
duality.  On the lisse Tate--Kummer subcategory it induces an involution with
\begin{equation}
 d_{\Uk}(Q)=Q^{-1},\qquad d_{\Uk}(\kappaell)=\kappaell,
 \qquad \kappaell^2=1.
 \label{eq:residual-coefficient-involution}
\end{equation}
We extend $d_{\Uk}$ coefficientwise to rational functions and let it fix the
grading variable $z$; inversion of $z$ is written separately.
The finite-type shell of radius $r$ has the class equality
\[
 [Rq_{r,!}\Qell]=[R\pi_{\mathbb T,!}\Qell]\Lell^{r-1}
 =(\Lell-\kappaell)\Lell^{r-1}
\]
in $K_{\Uk}$; no product isomorphism
$\mathscr Q_r\simeq\mathbb T\times\mathbb A^{r-1}$ is asserted.
The two terms of $\Lell-\kappaell$ are not a formal subtraction invented after
counting.  They are the two adjacent compactly supported cohomology groups of
the one-dimensional torus: the top Tate line and the degree-one orientation
line.

Assign the Weyl factor $Q^m$ and put $t=m-r$.  The corresponding residual
layer contributes
\begin{equation}
 Q^m[\mathscr Q_r]
 =Q^{m-r}(1-\kappaell Q)
 =Q^t-\kappaell Q^{t+1}.
 \label{eq:adjacent-shell-degrees}
\end{equation}
Define the compact-orbital coefficient
$I_m^{\mathrm{orb}}=1+\sum_{r=1}^m[\mathscr Q_r]$; by
Theorem~\ref{thm:PGL2-residual-orbit-comparison} this is the compact-shell
realization of the spherical orbital integral, and the first term is the
core.  Summing the layers gives
\begin{align}
 J_m^{\mathrm{orb}}:=Q^m I_m^{\mathrm{orb}}
 &=Q^m+\sum_{t=0}^{m-1}(Q^t-\kappaell Q^{t+1})
 \notag\\
 &=1+(1-\kappaell)\sum_{s=1}^{m}Q^s.
 \label{eq:shell-telescoping}
\end{align}
The coordinate
\begin{equation}
 \boxed{s=t+1=m-r+1}
 \label{eq:canonical-translation-residual}
\end{equation}
is therefore forced by the adjacent cohomological degrees in
\eqref{eq:adjacent-shell-degrees} and by telescoping of the shell filtration.
It is not a relabeling chosen to make a denominator palindromic; no analogous
formula for a general orbital family is inferred from this calculation.

Set $A=1-\kappaell$ and let
\[
 C=\{(m,s):m\ge0,\ 0\le s\le m\},
 \qquad F=\{(m,0):m\ge0\}.
\]
For an explicitly height-graded finite array whose $(m,s)$ component has been
written as $V_{m,s}(s)$, define its de-Tated height character by
\begin{equation}
 \operatorname{ch}_{\mathrm{ht}}(V)
 :=\sum_{m,s}[V_{m,s}]z^mu^s.
 \label{eq:height-character-definition}
\end{equation}
This operation is defined on the supplied height grading; it is not a functor
on the ungraded derived category and does not identify the twist $(s)$ with a
constant coefficient object.  Specialization $u\mapsto Q$ recovers the
ordinary Grothendieck class because $[V(s)]=Q^s[V]$ in the present convention.

Here $\one_C$ and $\one_F$ denote the constant closed-cone cellular
generators on the finite rational fan with maximal cone $C$ and face $F$;
they are not direct sums of one object for every lattice point.  Define the
auxiliary untwisted cellular envelope
\begin{equation}
 \widehat\Phi=A\,\one_C+\kappaell\,\one_F
 \in K_0\Cell_C(D_c^b(U,\Qell)).
 \label{eq:residual-skeleton-class}
\end{equation}
It encodes the de-Tated height character extracted from
\eqref{eq:shell-telescoping}; it is not asserted to be the ordinary
cellularization of the Tate-twisted shell complexes.  Its infinite lattice
expansion is interpreted only through the rational
generating transform in the localization generated by $1-z$ and $1-uz$.
Thus no infinite coproduct in a bounded cellular category is being asserted.
The bivariate lattice transform of $\widehat\Phi$ and its cellular partner are
\begin{equation}
 \mathcal W(z,u)
 =\frac{A}{(1-z)(1-uz)}+\frac{\kappaell}{1-z}
 =\frac{1-\kappaell uz}{(1-z)(1-uz)},
 \label{eq:residual-W}
\end{equation}
\begin{equation}
 \mathcal W^!(z,u)
 =A\frac{uz^2}{(1-z)(1-uz)}
 -\kappaell\frac{z}{1-z}.
 \label{eq:residual-partner}
\end{equation}
Theorem~\ref{thm:GF-duality}, with the displayed integral orientations, gives
\begin{equation}
 \mathcal W^!(z^{-1},u^{-1})=\mathcal W(z,u).
 \label{eq:residual-partner-inversion}
\end{equation}
Direct simplification, using $-\kappaell(1-\kappaell)=1-\kappaell$, also gives the
special self-partner identity
\begin{equation}
 \mathcal W^!(z,u)=-\kappaell z\,\mathcal W(z,u).
 \label{eq:residual-partner-identification}
\end{equation}
Specialize $u\mapsto Q$ and write
$W_{PGL_2}(z)=\mathcal W(z,Q)$.  Apply $d_{\Uk}$ to
the two preceding identities, using
\eqref{eq:residual-coefficient-involution}.  This yields the unambiguous
one-variable formula
\begin{equation}
 \boxed{W_{PGL_2}(z^{-1})
 =-\kappaell z\,d_{\Uk}(W_{PGL_2}(z)),
 \qquad \kappaell=[\mathfrak o_{\mathbb T}]=[\cK],\quad
 \mathfrak o_{\mathbb T}=\det(X_*(\mathbb T)\otimes\Qell).}
 \label{eq:quadratic-residual-orientation-reciprocity}
\end{equation}

Equation~\eqref{eq:quadratic-residual-orientation-reciprocity} is an
unconditional rational $K_0$ identity assembled from finite-type schemes.
Theorem~\ref{thm:PGL2-residual-orbit-comparison} identifies those schemes with
the compact positive-loop orbit shells that represent the actual spherical
$PGL_2$ double-coset masses.  Thus the equation is an orbital reciprocity
theorem for this family.  It is not a theorem about the entire derived quotient
stack of the affine-Springer locus.

\begin{proposition}[Interpretation of the three factors]
In~\eqref{eq:quadratic-residual-orientation-reciprocity}:
\begin{enumerate}[label=\textup{(\roman*)}]
\item $d_{\Uk}$ is the coefficient involution induced by dimension-normalized
  Verdier duality;
\item $\mathfrak o_{\mathbb T}$ is the determinant of the norm-one torus
  cocharacter local system and records quadratic descent;
\item $-z$ is the localized-$K_0$ shadow of the canonical module of the depth
  ray, including its odd shift and primitive translation.
\end{enumerate}
The coordinate $s=m-r+1$ belongs to the preceding orbital-shell-to-skeleton
bookkeeping step
and is not any one of these three factors.
\end{proposition}

\begin{proof}
Part (i) is the definition in
\eqref{eq:residual-coefficient-involution}.  Part (ii) is the quadratic sign
action on $X_*(\mathbb T)$, together with
Theorem~\ref{thm:torus-Frobenius-bimodule}; it uses no orbital product chart.
Part (iii) is the canonical module of $\mathbf N$.
Equation~\eqref{eq:canonical-translation-residual} was derived before applying
polyhedral duality, so it belongs to the orbital shell grading convention, not to the same-series
identification~\eqref{eq:residual-partner-identification}.
\end{proof}

\begin{warning}[The functional equation is still a $K_0$ statement]
The class $A=1-\kappaell$ is virtual, and the denominators in
\eqref{eq:residual-W} represent an infinite grading.  The equality
$-\kappaell(1-\kappaell)=1-\kappaell$ uses cohomological signs in $K_0$ and is not an
isomorphism between the naive bounded complexes representing the two sides.
The finite labelled-shell theorem below gives an object-level partner
equivalence on its discrete target.  The more polyhedral height object exists
only after cohomological associated graded, and its assembled horizontal
duality is asserted only after $K_0$.  Neither theorem constructs cross-shell
transition maps or an incidence-bearing ungraded equivalence.  Completion is
a separate operation even for an eventually constant finite filtration:
$t$-adic completion generally destroys the $t=1$ specialization.  An infinite
filtration adds further continuity and convergence questions.
\end{warning}


\section{A finite labelled-shell Verdier partner and its polyhedral
\texorpdfstring{$K_0$}{K0} shadow}
\label{sec:PGL2-narrow-skeletonization}

This section separates two statements of different strength.  First, the
actual finite family of labelled orbit shells admits a genuine
zero-dimensional Verdier partner: both direct-image functors are defined
independently and Verdier duality intertwines them by a natural equivalence.
Second, the height regrading produces a more polyhedral object only after
cohomological associated graded; its assembled horizontal duality remains a
$K_0$ statement.

\subsection{The finite labelled-shell Verdier partner}

Fix $N\ge0$ and let
\begin{equation}
 \Lambda_N=\{(m,r):0\le r\le m\le N\}
 \label{eq:PGL2-finite-shell-index-set}
\end{equation}
with the discrete finite cell structure, every point being a zero-cell.  Put
\begin{equation}
 \mathscr Y_N^{\mathrm{orb}}
 =\coprod_{(m,r)\in\Lambda_N}\mathscr O_{m,r},
 \qquad
 \mathscr Z_N=U\times\Lambda_N
 =\coprod_{(m,r)\in\Lambda_N}U,
 \label{eq:PGL2-finite-shell-source-target}
\end{equation}
Write $q_{m,r}:\mathscr O_{m,r}\to U$ for the structure map and define the
labelled shell map
\begin{equation}
 \rho_N:\mathscr Y_N^{\mathrm{orb}}\longrightarrow\mathscr Z_N,
 \qquad
 x\in\mathscr O_{m,r}\longmapsto(q_{m,r}(x),(m,r)).
 \label{eq:PGL2-labelled-shell-map}
\end{equation}
This map is constructed from the actual orbit schemes and their structure
maps; the target is not inferred from a generating function.

There is a canonical equivalence
\begin{equation}
 D_c^b(\mathscr Z_N,\Qell)
 \simeq
 \Cell_{\Lambda_N}(D_c^b(U,\Qell))
 \simeq
 \prod_{(m,r)\in\Lambda_N}D_c^b(U,\Qell).
 \label{eq:PGL2-zero-cellular-target}
\end{equation}
Let $\D_{\mathrm{cell},N}$ denote total cellular duality on this category.
Because all cells have dimension zero, it is componentwise Verdier duality
$\D_U$; there is no additional cellular shift or orientation line.  Define
the two functors independently by
\begin{equation}
 \Sk_{N,!}=R\rho_{N,!},\qquad
 \Sk_{N,*}=R\rho_{N,*}:
 D_c^b(\mathscr Y_N^{\mathrm{orb}},\Qell)
 \longrightarrow
 \Cell_{\Lambda_N}(D_c^b(U,\Qell)),
 \label{eq:PGL2-genuine-Sk-pair}
\end{equation}
using~\eqref{eq:PGL2-zero-cellular-target}.  In particular $\Sk_{N,*}$ is
not defined as a Verdier conjugate of $\Sk_{N,!}$.

\begin{theorem}[Finite labelled-shell Verdier partner]
\label{thm:PGL2-genuine-small-skeletonization}
For every finite cutoff $N$, the functors in
\eqref{eq:PGL2-genuine-Sk-pair} are defined on the full bounded constructible
category and there is a canonical natural equivalence of contravariant
functors
\begin{equation}
 \boxed{
 \D_{\mathrm{cell},N}\circ\Sk_{N,!}
 \simeq
 \Sk_{N,*}\circ\D_{\mathscr Y_N^{\mathrm{orb}}}.}
 \label{eq:PGL2-genuine-small-skeletonization}
\end{equation}
For $M=(M_{m,r})$ under the finite disjoint-union decomposition, its
$(m,r)$ component is precisely
\begin{equation}
 \D_U Rq_{m,r,!}M_{m,r}
 \simeq
 Rq_{m,r,*}\D_{\mathscr O_{m,r}}M_{m,r}.
 \label{eq:PGL2-genuine-component-partner}
\end{equation}
The equivalences are compatible with restriction from $N+1$ to $N$.
\end{theorem}

\begin{proof}
Each $\mathscr O_{m,r}\simeq\mathscr Q_r$ is a smooth affine finite-type
$U$-scheme by Theorems~\ref{thm:finite-residual-filtration} and
\ref{thm:PGL2-residual-orbit-comparison}.  Hence $\rho_N$ is separated and of
finite type, and both direct images preserve bounded constructibility.  The
six-functor Verdier isomorphism for this independently specified map gives
\[
 \D_{\mathscr Z_N}R\rho_{N,!}
 \simeq
 R\rho_{N,*}\D_{\mathscr Y_N^{\mathrm{orb}}}
\]
naturally on the full source category.  Under
\eqref{eq:PGL2-zero-cellular-target}, $\D_{\mathscr Z_N}$ is exactly
$\D_{\mathrm{cell},N}$, proving the boxed equivalence.  Decomposition over the
zero-cells gives~\eqref{eq:PGL2-genuine-component-partner}.  The inclusion
$\Lambda_N\subset\Lambda_{N+1}$ is a union of open-and-closed components, so
restriction preserves the construction and its Verdier transformation.
\end{proof}

\begin{remark}[Why the theorem is genuine and why it is small]
The source, target, map, two functors, and both dualities are all fixed before
the equivalence is invoked, and the equivalence is natural in every
constructible coefficient object.  Thus
Theorem~\ref{thm:PGL2-genuine-small-skeletonization} is an object-level
Verdier partner identity, not a $K_0$ identity.  It is small because
$\Lambda_N$ is discrete: the theorem records the actual labels $(m,r)$ but no
closure incidence, no height translation $s=m-r+1$, and no Stanley
canonical-module operation.  The split product chart is not needed for this
theorem.
\end{remark}

\subsection{The associated-graded polyhedral shadow}

We next retain the canonical height data in addition to the shell labels.  Let
\[
 C=\{(m,s)\in\mathbf N^2:0\le s\le m\},
 \qquad F=\{(m,0):m\ge0\}.
\]
For a complex $M$, write
$\operatorname{gr}_{\mathcal H}M=\bigoplus_i\mathcal H^i(M)[-i]$.
This associated graded is canonical; it does not assert that $M$ splits.

Retaining the structure maps $q_{m,r}$ above, put
\begin{equation}
 B_{m,0}=\Qell(m),\qquad
 B_{m,r}=Rq_{m,r,!}\Qell(m)[2r]\quad(1\le r\le m).
 \label{eq:PGL2-normalized-shell-complexes}
\end{equation}
The shift only places the two compact-cohomology sheaves in degrees $-1$ and
$0$; its class is unaffected.

\begin{theorem}[Associated-graded coefficient array and its $K_0$ polyhedral shadow]
\label{thm:PGL2-narrow-skeletonization}
For the actual compact orbit shells of
Theorem~\ref{thm:PGL2-residual-orbit-comparison}, the following statements
hold.
\begin{enumerate}[label=\textup{(\roman*)}]
\item For $r\ge1$ there is a canonical identification of cohomology sheaves
  \begin{equation}
   \operatorname{gr}_{\mathcal H}B_{m,r}
   \simeq
   \Qell(m-r)\oplus\cK(m-r+1)[1].
   \label{eq:PGL2-shell-cohomology-graded}
  \end{equation}
  Consequently the canonical translation $s=m-r+1$ gives
  \begin{equation}
   \bigoplus_{r=0}^m\operatorname{gr}_{\mathcal H}B_{m,r}
   \simeq
   \Qell\oplus
   \bigoplus_{s=1}^m(\Qell\oplus\cK[1])(s).
   \label{eq:PGL2-canonical-reindexing-object}
  \end{equation}
\item For a finite cutoff $N$, put
  $C_N=\{(m,s)\in C:m\le N\}$ and define the discrete graded-array category
  \begin{equation}
   \operatorname{Gr}_{C_N}(U)
   :=\prod_{(m,s)\in C_N}D_c^b(U,\Qell).
   \label{eq:PGL2-discrete-graded-array-category}
  \end{equation}
  The component formula
  \begin{equation}
   \Psi^{\operatorname{gr}}_{N;m,s}
   =
   \begin{cases}
    \Qell,&s=0,\\
    (\Qell\oplus\cK[1])(s),&1\le s\le m
   \end{cases}
   \label{eq:PGL2-discrete-graded-array}
  \end{equation}
  defines an object of~\eqref{eq:PGL2-discrete-graded-array-category}, and
  \eqref{eq:PGL2-canonical-reindexing-object} identifies it with the
  cohomological associated graded of the normalized actual shell coefficients.
  This is a discrete coefficient array, not yet a cellular sheaf on the
  incidence-bearing cone $C$.
\item Apply the de-Tated height character
  \eqref{eq:height-character-definition} to
  \eqref{eq:PGL2-discrete-graded-array}.  The resulting class function is the
  restriction to $C_N$ of the auxiliary untwisted cellular envelope
  \begin{equation}
   \widehat\Phi=(1-\kappaell)\one_C+\kappaell\one_F
   \in K_0\Cell_C(D_c^b(U,\Qell)),
   \label{eq:PGL2-associated-graded-K0-envelope}
  \end{equation}
  from~\eqref{eq:residual-skeleton-class}.  This is not an equality between
  the Tate-twisted shell array and an ordinary cellular object: the twist is
  recorded by $u^s$, and $u\mapsto Q$ recovers its ordinary $K_0$ class.
  Nor is $\widehat\Phi$ the image of a constructed geometric functor to $C$.
  Cellular duality and the lattice transform of this auxiliary envelope give
  \eqref{eq:PGL2-actual-orbital-reciprocity} after rational continuation and
  the specialization $u\mapsto Q$.
\end{enumerate}
\end{theorem}

\begin{proof}
The comparison theorem reduces the compact cohomology to $H_r$.  The two
nonzero cohomology sheaves of $R\pi_{\mathbb T,!}\Qell$ are
\[
 \mathcal H^1=\cK,
 \qquad \mathcal H^2=\Qell(-1).
\]
The vector-torsor Thom shift in~\eqref{eq:Haar-shell-Thom} therefore gives
\[
 \mathcal H^{2r-1}(Rq_{m,r,!}\Qell)=\cK(-(r-1)),
 \qquad
 \mathcal H^{2r}(Rq_{m,r,!}\Qell)=\Qell(-r).
\]
Twisting by $(m)$ and shifting by $[2r]$ proves
\eqref{eq:PGL2-shell-cohomology-graded}.  Its Tate terms have heights
$m-r=0,\ldots,m-1$, the core supplies height $m$, and its Kummer terms have
heights $m-r+1=1,\ldots,m$.  This proves
\eqref{eq:PGL2-canonical-reindexing-object} and
\eqref{eq:PGL2-discrete-graded-array}.

After removing the displayed twist $(s)$ into the formal height monomial
$u^s$, the class of $\Qell\oplus\cK[1]$ is $1-\kappaell$, while the class at
$s=0$ is $1=(1-\kappaell)+\kappaell$.  This proves the height-character statement
\eqref{eq:PGL2-associated-graded-K0-envelope}; it would be false as a literal
identification of the Tate-twisted array with an untwisted ordinary cellular
object.  The calculation in
Section~\ref{sec:quadratic-residual-series}, equivalently
Theorem~\ref{thm:GF-duality} applied to the auxiliary envelope and then
specialized by $u\mapsto Q$, gives the functional equation.
\end{proof}

\begin{remark}[The split product chart computes the vertical partner only]
After the Kummer cover $b^2=a$, the ratio of the two split factors gives an
explicit chart
\begin{equation}
 \mathscr Q_r\times_U\widetilde U
 \simeq
 \operatorname{Res}_{\widetilde U[\varpi]/(\varpi^r)/\widetilde U}
 \mathbb G_m
 \simeq\mathbb G_m\times\mathbb A^{r-1}.
 \label{eq:PGL2-split-tame-product-chart}
\end{equation}
Thus compact and ordinary push-forwards are independently computed by the
product chart.  Fixed-map Verdier duality
$\D_U Rq_!\simeq Rq_*\D$ and
Theorem~\ref{thm:torus-Frobenius-bimodule} identify their vertical coefficient
partners.  The deck involution interchanges the two
quadratic factors; it acts by $-1$ on the rank-one torus cocharacter and hence
descends as $\cK$.  All product-chart pairings are equivariant, so the
shellwise object comparison descends to $U$.  The chart supplies no geometric
map to $C$ and no incidence morphisms between distinct $(m,r)$.
\end{remark}

\begin{warning}[The zero-dimensional theorem does not prove the cone theorem]
The source and target of
Theorem~\ref{thm:PGL2-genuine-small-skeletonization} are fixed independently,
so its partner equivalence is genuine.  By contrast,
Theorem~\ref{thm:PGL2-narrow-skeletonization} deliberately passes to
cohomological associated graded and then to a cellular envelope only after
$K_0$.  It defines no geometric skeleton map to $C$, no functor on arbitrary
orbital morphisms, and no cross-shell incidence maps.  It therefore proves
neither an object-level positive-dimensional identity
$\Dcell\Sk_!\simeq\Sk_*\D$, nor a Rees action or compatibility with
convolution.
\end{warning}



\appendix
\section{Cellular duality and the generating-function shadow}

Let $N$ be a lattice and let $\Sigma$ be a finite pointed rational fan in
$N_\mathbf R$.  For a cell $\sigma$, let $\mathbf k_{\bar\sigma}$ denote the
constant sheaf on its closure, extended by zero, and let
$\mathbf k_{\sigma^\circ}$ denote extension by zero from its relative
interior.  Scalar cellular duality gives
\begin{equation}
 \Dcell\mathbf k_{\bar\sigma}
 \simeq
 \operatorname{or}_\sigma[\dim\sigma]
 \otimes\mathbf k_{\sigma^\circ}.
 \label{eq:cell-generator-duality}
\end{equation}
The link complex, rather than a rank-one interior indicator, is used at a
singular face.  Cellular sheaf--cosheaf duality and subdivision invariance are
developed in \cite{Curry}; duality and direct image for sheaves on fans are
also treated in \cite{BresslerLunts}.

Let $\cC$ be a coefficient category with a duality $d$.  We use
$\Cell_\Sigma(\cC)$ for $\cC$-valued cellular diagrams and write
$\D^{\mathrm{tot}}_{\cC,\Sigma}=d\boxtimes\Dcell$.  Thus
\[
 \D^{\mathrm{tot}}_{\cC,\Sigma}
 (V\boxtimes\mathbf k_{\bar\sigma})
 \simeq
 d(V)\otimes\operatorname{or}_\sigma[\dim\sigma]
 \otimes\mathbf k_{\sigma^\circ}.
\]
The orientation line in~\eqref{eq:cell-generator-duality} is part of
$\Dcell$; it must not be
silently replaced by a scalar sign when the integral lattice has monodromy.
Define the coefficient--orientation involution on closed-cell generators by
\begin{equation}
 \mathfrak d^{\operatorname{or}}_\Sigma
 \bigl(V\otimes\mathbf k_{\bar\sigma}\bigr)
 =d(V)\otimes\operatorname{or}_\sigma
  \otimes\mathbf k_{\bar\sigma}.
 \label{eq:coefficient-orientation-duality}
\end{equation}
Unlike $\D^{\mathrm{tot}}_{\cC,\Sigma}$, this operation does not exchange a closed cell with
its interior.

For $\Phi\in K_0\Cell_\Sigma(\cC)$, its stalk class is a cellwise-constant
function
\[
 \varphi_\Phi:N\cap|\Sigma|\longrightarrow K_0(\cC).
\]
When the resulting series is rational, define
\begin{equation}
 \GF(\Phi)(z)=
 \sum_{u\in N\cap|\Sigma|}\varphi_\Phi(u)z^u
 \label{eq:GF-definition}
\end{equation}
in the localization generated by the relevant factors $1-z^v$.

\begin{theorem}[Generating functions transport cellular duality]
\label{thm:GF-duality}
For every finite $K_0$-combination of closed rational cones with coefficients
in $\cC$,
\begin{equation}
 \boxed{
 \GF\bigl(\D^{\mathrm{tot}}_{\cC,\Sigma}\Phi\bigr)(z^{-1})
 =\GF\bigl(\mathfrak d^{\operatorname{or}}_\Sigma\Phi\bigr)(z).}
 \label{eq:GF-duality}
\end{equation}
The equality is independent of rational polyhedral subdivision.
If the integral orientations are compatibly trivialized, the right side
reduces to $d(\GF(\Phi)(z))$.
\end{theorem}

\begin{proof}
The classes $[V\otimes\mathbf k_{\bar\sigma}]$ generate the cellular
Grothendieck group.  Put $e=\dim\sigma$.  On this generator the left side is
\[
 (-1)^e[d(V)\otimes\operatorname{or}_\sigma]
 \sum_{u\in\relint(\sigma)\cap N}z^{-u}.
\]
Rational cone reciprocity \cite{Stanley} gives
\[
 \sum_{u\in\relint(\sigma)\cap N}z^{-u}
 =(-1)^e\sum_{u\in\sigma\cap N}z^u.
\]
The two signs cancel, proving~\eqref{eq:GF-duality} on generators and hence
by additivity in general.  Under a subdivision, inclusion--exclusion on
closed cells and its Verdier dual include all internal faces with the
correct link multiplicity.  Equivalently, the constructible stalk function
and the dualizing complex do not depend on the chosen cellular presentation.
\end{proof}


\section{The local unramified absolute-\texorpdfstring{$A_1$}{A1} theorem}
\label{sec:absolute-A1}

The represented $PGL_2$ family separates the Kummer line from the Tate
coefficient.  We now test whether that line survives changes of central
isogeny and shell lattice.  This section is local: except in the saturated
$PGL_2$ case, its coefficient objects are explicit orbit-summed
interpolation classes, not relative orbital push-forwards.

Let $S/k$ be a smooth connected cameral parameter base of dimension $d$ over
a finite field $k$, with $\ell\ne\operatorname{char}k$.  Write $\Lell$ also
for the class $[\Qell(-1)[-2]]$ over $S$, and put
\[
 \mathcal R_S
 =K_0\bigl(D^b_{\mathrm{lisse,Weil}}(S,\Qell)\bigr)
  \otimes_{\mathbf Z}\mathbf Q,
 \qquad Q=\Lell^{-1}.
\]
Normalize coefficient duality by
\begin{equation}
 d_S(M)=\D_S(M)(-d)[-2d].
 \label{eq:A1-normalized-duality}
\end{equation}
Thus $d_S$ is contragredience on lisse objects and sends $Q$ to $Q^{-1}$.
All rational identities below lie in the $d_S$-stable localization
\begin{equation}
 \Lambda_S
 =\mathcal R_S[z^{\pm1},
 (1-z)^{-1},(1+z)^{-1},
 (1-Qz)^{-1},(1-Q^{-1}z)^{-1}].
 \label{eq:A1-localization}
\end{equation}
Let $\widetilde S\to S$ be the quadratic cameral cover of the root torus.
Its sign line is the rank-one local system $\mathfrak o$; at a closed point
$x$ write $\chi_x\in\{1,-1\}$ for its geometric-Frobenius trace.  The deck
involution acts on the root cocharacter line by $-1$, so throughout this
section
\begin{equation}
 \mathfrak o\simeq\mathfrak o_{G,T}
 =\det\!\left(X_*(T/Z(G)^\circ)\otimes\Qell\right).
 \label{eq:A1-cameral-orientation}
\end{equation}

\begin{remark}[Central neutrality is necessary]
\label{rem:central-neutral-necessary}
For $G=GL_2$, multiplying a standard lift by $\varpi I$ preserves its image
in $PGL_2$, quadratic order, fixed set, Weyl discriminant, and root
orientation.  Its determinant then has nonzero valuation, so no conjugate
lies in $GL_2(\cO_F)$ and the spherical orbital integral vanishes.  Condition
\eqref{eq:central-neutral-support} removes exactly this connected-center
counterexample; alternatively one would have to change the test function to
one compact modulo center.
\end{remark}

Let $p:G\to G_{\ad}$, put $H=p(G(F))$, and let $\mathscr B$ be the reduced
Bruhat--Tits tree of $PGL_2(F)$.

\begin{lemma}[Finite central lifting over the hyperspecial model]
\label{lem:A1-central-lifting}
Put $G'=G/Z(G)^\circ$ and let $\mathcal Z^\circ$ be the smooth connected
hyperspecial model of the unramified torus $Z(G)^\circ$.  After removing the connected center, the induced
central isogeny has the form
\begin{equation}
 1\longrightarrow Z_{\mathrm{fin}}
 \longrightarrow G'\longrightarrow G_{\ad}\longrightarrow1,
 \qquad |Z_{\mathrm{fin}}|\le2.
 \label{eq:A1-finite-central-isogeny}
\end{equation}
For the hyperspecial models one has
\[
 H^1(\cO_F,\mathcal Z^\circ)=1,
 \qquad
 H^1(\cO_F,Z_{\mathrm{fin}})\hookrightarrow H^1(F,Z_{\mathrm{fin}}).
\]
Consequently an element of $G_{\ad}(\cO_F)$ which has a lift to $G(F)$
already lifts to $K=\mathcal G(\cO_F)$.
\end{lemma}

\begin{proof}
The connected center is an unramified torus.  Its smooth integral model has
connected special fiber, so Lang's theorem followed by henselian lifting
gives $H^1(\cO_F,\mathcal Z^\circ)=1$.  The remaining kernel is a quotient of
the center of the simply connected group of type $A_1$, hence has order at
most two.  In odd residue characteristic it is finite \'etale.  A finite
\'etale torsor over the henselian discrete valuation ring which is trivial on
the generic fiber is already trivial over the ring, proving the displayed
injectivity.  The final assertion is the connecting-class criterion applied
first to the finite central isogeny and then to the connected center.
\end{proof}

\begin{lemma}[Classification reduction for the reduced-building data]
\label{lem:A1-building-classification}
For every pair and family in Definitions~\ref{def:A1-pair} and
\ref{def:standard-A1-family}:
\begin{enumerate}[label=\textup{(\alph*)}]
\item $G_{\ad}\simeq PGL_2$, $G_{\der}^{\mathrm{sc}}\simeq SL_2$, and
  multiplication induces a central isogeny
  $Z(G)^\circ\times SL_2\to G$.
\item The type character
  \begin{equation}
   \tau:PGL_2(F)\longrightarrow\mathbf Z/2,
   \qquad[g]\longmapsto\ord_F(\det g)\pmod2,
   \label{eq:A1-type-character}
  \end{equation}
  has $\tau(H)=0$ or $\mathbf Z/2$.  In the first case $H$ preserves vertex
  color; in the second it sees both colors.
\item The root torus $T_{\rt}$ is split or the norm-one torus of the
  unramified quadratic extension.  Its orientation line is, respectively,
  the constant line or the unramified quadratic sign line.
\item If $K_{\ad}=G_{\ad}(\cO_F)$, then
  \begin{equation}
   p(K)=H\cap K_{\ad},
   \qquad
   T(F)\backslash G(F)\xrightarrow{\sim}
   (H\cap T_{\rt}(F))\backslash H.
   \label{eq:A1-adjoint-quotient}
  \end{equation}
\item The shell lattice is saturated when $\tau(H)=\mathbf Z/2$ and
  type-preserving when $\tau(H)=0$.  In the type-preserving elliptic case
  there are exactly two rational embedding phases, according to the color of
  the fixed vertex; they retain the even or the odd spheres.  This list is
  exhaustive under the two definitions above.
\end{enumerate}
\end{lemma}

\begin{proof}
The $A_1$ Dynkin diagram has no nontrivial automorphism.  Hence the
quasi-split adjoint form is split $PGL_2$, its simply connected cover is
$SL_2$, and the center--derived-group decomposition proves (a)
\cite[Proposition~20.32 and Theorem~23.51]{MilneAlgebraicGroups}; see also
\cite[Appendix~C.4]{ConradReductiveGroups}.

The character~\eqref{eq:A1-type-character} is well defined because scalar
multiplication changes determinant valuation by an even integer.  The image
$H$ contains $PSL_2(F)$, which is transitive on the vertices of either fixed
color.  This proves (b).  Unramified maximal tori in $PGL_2$ come from the
split quadratic algebra or the unramified quadratic field, and Galois acts
by $+1$ or $-1$ on their cocharacter lines, proving (c).

The kernel of $G(F)\to H$ lies in $T(F)$, which gives the quotient
homeomorphism in (d).  Lemma~\ref{lem:A1-central-lifting} shows that an
integral adjoint point with a generic lift already lifts integrally; hence
$p(K)=H\cap K_{\ad}$.  Finally the split torus fixes an apartment containing
both colors, whereas an unramified elliptic torus fixes one vertex.
Intersecting the standard tube or ball with the $H$-orbit retains all
spheres, the even spheres, or the odd spheres.  A type-changing element
exchanges the two elliptic phases, and $PSL_2(F)$ is transitive within each
color.  This proves (e) and exhaustiveness.
\end{proof}

The preceding lemma classifies only reduced-building data.  It does not
identify the number or volume of the $T(F)$-orbits inside a sphere; that is
the content of the next lemma.

\begin{lemma}[Adjoint quotient measure and orbit-summed shell transfer]
\label{lem:A1-shell-transfer}
Put
\[
 T_H=H\cap T_{\rt}(F),\qquad K_H=H\cap K_{\ad}.
\]
Transport the quotient measure through~\eqref{eq:A1-adjoint-quotient} and
give $K$, $T(F)\cap K$, $K_H$, and $T_H\cap K_H$ volume one.  If the root
torus has Frobenius sign $\chi\in\{1,-1\}$, every allowed positive sphere of
radius $r$ has total mass
\begin{equation}
 \sum_{\mathcal O\text{ over the radius-}r\text{ sphere}}
 \vol(T_H\cap g_{\mathcal O}K_Hg_{\mathcal O}^{-1})^{-1}
 =(q-\chi)q^{r-1}.
 \label{eq:A1-orbit-summed-shell}
\end{equation}
The radius-zero core, when allowed, has mass one.  Every radius is allowed in
the saturated case.  In a type-preserving elliptic phase $\epsilon$, exactly
the radii $r\equiv\epsilon\pmod2$ are allowed; in the split case every
transverse radius is allowed.
\end{lemma}

\begin{proof}
The two sides of~\eqref{eq:A1-adjoint-quotient} carry Haar quotient measures,
so their ratio is constant.  The compact core has mass one on both sides by
Weil's formula and $p(K)=K_H$, hence the ratio is one.

In the elliptic case $T_H$ is compact.  Orbit--stabilizer gives
\[
 \sum_{\mathcal O\subset\Sigma_r}
 \vol(T_H\cap g_{\mathcal O}K_Hg_{\mathcal O}^{-1})^{-1}
 =|\Sigma_r|=(q+1)q^{r-1}.
\]
The sum is essential: for example, the first $SL_2$ elliptic sphere splits
into two norm-one-torus orbits even though their total mass remains $q+1$.
In the split case choose a fundamental domain for the translation lattice
along the apartment.  Its transverse radius-$r$ section has
$(q-1)q^{r-1}$ vertices.  If $\tau(H)=\mathbf Z/2$, choose
$h=[g]\in H$ with $\ord_F(\det g)$ odd and put $a=\det g$.  The kernel of
the determinant square-class map
$PGL_2(F)\to F^\times/(F^\times)^2$ is $PSL_2(F)$; hence $h$ differs from
$[\operatorname{diag}(a,1)]$ by an element of $PSL_2(F)\subset H$.  The
latter diagonal element therefore lies in $H$ and translates the chosen
apartment by odd length.  If $\tau(H)=0$, only the longitudinal period
changes.  In either case orbit--stabilizer over a longitudinal fundamental
domain leaves the transverse orbit-summed mass unchanged.
The allowed parities are those of
Lemma~\ref{lem:A1-building-classification}, proving the result.
\end{proof}

Combining Lemma~\ref{lem:A1-shell-transfer} with
\eqref{eq:central-neutral-support} removes the connected center from the
spherical orbital integral.  For a closed point $x\to S$, replace
$(\Lell,\mathfrak o)$ by $(q_x,\chi_x)$ in the following displayed shell
classes.  Their Frobenius traces are the corresponding local orbital masses.
In the saturated $PGL_2$ case they are the genuine residual push-forwards of
Theorem A; otherwise they are the announced orbit-summed interpolation.

\subsection{The \texorpdfstring{$SL_2$}{SL2} parity test}

The reduced buildings of $SL_2$ and $PGL_2$ agree, but $SL_2(F)$ preserves
vertex color.  This changes the elliptic shell semigroup without changing
the root orientation.

\begin{proposition}[Type-preserving shell coefficients]
\label{prop:SL2-parity-coefficients}
At a split point,
\begin{equation}
 I_{+,m}=\Lell^m,\qquad W_+(z)=\frac1{1-z}.
 \label{eq:SL2-split-series}
\end{equation}
At an unramified elliptic point let $\epsilon=0$ when the fixed vertex has
the hyperspecial color and $\epsilon=1$ otherwise.  If $m=2e$ or $2e+1$,
then phase $0$ has
\begin{align}
 I_{-,2e}^{(0)}=I_{-,2e+1}^{(0)}
 &=\sum_{r=0}^{2e}\Lell^r,\notag\\
 J_{-,2e}^{(0)}&=\sum_{s=0}^{2e}Q^s,
 &J_{-,2e+1}^{(0)}&=\sum_{s=1}^{2e+1}Q^s,
 \label{eq:SL2-even-phase}
\end{align}
whereas phase $1$ has
\begin{align}
 I_{-,2e}^{(1)}&=\sum_{r=0}^{2e-1}\Lell^r,
 &J_{-,2e}^{(1)}&=\sum_{s=1}^{2e}Q^s,\notag\\
 I_{-,2e+1}^{(1)}&=\sum_{r=0}^{2e+1}\Lell^r,
 &J_{-,2e+1}^{(1)}&=\sum_{s=0}^{2e+1}Q^s.
 \label{eq:SL2-odd-phase}
\end{align}
The empty sum gives $I_{-,0}^{(1)}=0$.
\end{proposition}

\begin{proof}
In the split case, passing to one vertex color changes translations along
the apartment but removes no transverse distance.  Lemma
\ref{lem:A1-shell-transfer} gives total positive shell mass
$(\Lell-1)\Lell^{r-1}$, and adding the core yields $\Lell^m$.
In elliptic phase $0$ the ball meets the chosen color in the even spheres,
including the core; in phase $1$ it meets it in the odd spheres and excludes
the core.  Summing the masses
$(\Lell+1)\Lell^{r-1}$ over the allowed radii gives the displayed $I$'s.
Multiplication by $Q^m$ gives the $J$'s.
\end{proof}

Put
\[
 e_+=\frac{1+\mathfrak o_{G,T}}2,
 \qquad e_-=\frac{1-\mathfrak o_{G,T}}2.
\]
The two elliptic generating functions and their split/elliptic interpolation
are
\begin{align}
 W_-^{(0)}(z)&=\frac{1+Qz^2}{(1-z^2)(1-Qz)},
 \label{eq:SL2-even-series}\\
 W_-^{(1)}(z)&=\frac{(1+Q)z}{(1-z^2)(1-Qz)},
 \label{eq:SL2-odd-series}\\
 W_{\mathrm{tp}}^{(\epsilon)}(z)
 &=e_+W_+(z)+e_-W_-^{(\epsilon)}(z).
 \label{eq:SL2-type-preserving-series}
\end{align}

\begin{proposition}[Parity changes the shell lattice, not the orientation]
\label{prop:SL2-parity-reciprocity}
For both elliptic embedding phases,
\begin{equation}
 W_{\mathrm{tp}}^{(\epsilon)}(z^{-1})
 =-z\,(\mathfrak o_{G,T}\otimes d_S)
 W_{\mathrm{tp}}^{(\epsilon)}(z).
 \label{eq:SL2-parity-reciprocity}
\end{equation}
\end{proposition}

\begin{proof}
One has $W_+(z^{-1})=-zW_+(z)$.  Direct inversion of
\eqref{eq:SL2-even-series} and~\eqref{eq:SL2-odd-series}, using
$d_S(Q)=Q^{-1}$, gives
$W_-^{(\epsilon)}(z^{-1})=z\,d_SW_-^{(\epsilon)}(z)$.  The orientation line
has eigenvalue $+1$ on the split projector and $-1$ on the elliptic
projector, so the identities combine as claimed.
\end{proof}

The saturated $PGL_2$ series and the two $SL_2$ phase series are distinct
rational functions supported on different shell semigroups.  Their common
inversion line is therefore not determined by the shell lattice.

\subsection{The \texorpdfstring{$U_2$}{U2} connected-center test}

\begin{proposition}[The connected center forces root orientation]
\label{prop:U2-center-test}
Let $E/F$ be unramified quadratic and let $U_2$ be the quasi-split unitary
group of an isotropic hermitian plane.  Then
\[
 U_2^{\der}=SU_2\simeq SL_2,\qquad U_2^{\ad}\simeq PGL_2,
\]
and its norm-one connected center acts trivially on the reduced building.
For central-neutral spherical families, reciprocity is controlled by
\[
 \mathfrak o_{U_2,T}
 =\det\!\left(X_*(T/Z(U_2)^\circ)\otimes\Qell\right),
\]
whereas replacing this line by $\det(X_*(T)\otimes\Qell)$ gives a false
identity in the split-root specialization.
\end{proposition}

\begin{proof}
The central isogeny
$Z(U_2)^\circ\times SU_2\to U_2$ is the rank-two instance of the
center--derived-group decomposition.  The adjoint image is type-preserving;
equivalently, the connecting map for
$1\to U_1\to U_2\to PGL_2\to1$ lands in
$H^1(F,U_1)=F^\times/\Nm(E^\times)$ and detects valuation parity.  Put
$\mathfrak o_Z=\det(X_*(Z(U_2)^\circ)\otimes\Qell)$.  The rational exact
sequence of cocharacter spaces gives
\[
 \det(X_*(T)\otimes\Qell)
 =\mathfrak o_Z\otimes\mathfrak o_{U_2,T}.
\]
For a central-neutral split-root family the adjoint series is
$W_+(z)=1/(1-z)$ and satisfies $W_+(z^{-1})=-zW_+(z)$.  The nontrivial
central sign $\mathfrak o_Z$ is invisible to both support and building
geometry, so inserting it would contradict this identity.  Hence the
connected-center direction must be deleted.
\end{proof}

\subsection{Local classification and the unified theorem}

The saturated interpolation, with $\mathfrak o=\mathfrak o_{G,T}$, is
\begin{equation}
 W_{\mathrm{sat}}(z)
 =\frac{1-\mathfrak o Qz}{(1-z)(1-Qz)},
 \qquad
 W_{\mathrm{sat}}(z^{-1})
 =-z\,(\mathfrak o\otimes d_S)W_{\mathrm{sat}}(z).
 \label{eq:A1-saturated-reciprocity}
\end{equation}

\begin{theorem}[Local orbital shell classification]
\label{thm:A1-local-shell-classification}
Let $(G,K)$ and $\gamma_m$ satisfy Definitions~\ref{def:A1-pair} and
\ref{def:standard-A1-family}.  The Haar orbital coefficients of $\one_K$ are
the Frobenius specializations of either
\begin{equation}
 I_{G,m}^{\mathrm{sat}}
 =1+\sum_{r=1}^m(\Lell-\mathfrak o_{G,T})\Lell^{r-1},
 \label{eq:A1-saturated-shell-class}
\end{equation}
or the type-preserving formulas in
Proposition~\ref{prop:SL2-parity-coefficients}, according to the type image
and elliptic phase.  This is an equality of local orbit-summed masses.  It is
a relative orbital push-forward only in the represented saturated $PGL_2$
family of Theorem A.
\end{theorem}

\begin{proof}
Central neutrality pulls the support condition back from the adjoint group.
Lemma~\ref{lem:A1-building-classification} gives the exhaustive allowed
radii and phases, while Lemma~\ref{lem:A1-shell-transfer} gives the core and
every positive orbit-summed shell mass.  Summing them proves the formulas and
their Frobenius specialization.
\end{proof}

\begin{theorem}[Unramified rank-one orientation reciprocity]
\label{thm:A1-orientation-reciprocity}
Under the hypotheses of Theorem~\ref{thm:A1-local-shell-classification}, let
$W_G$ be the saturated or type-preserving interpolation selected by the type
image and elliptic phase.  In $\Lambda_S$ one has
\begin{equation}
 \boxed{
 W_G(z^{-1})
 =-z\,(\mathfrak o_{G,T}\otimes d_S)W_G(z),
 \qquad
 \mathfrak o_{G,T}
 =\det\!\left(X_*(T/Z(G)^\circ)\otimes\Qell\right).}
 \label{eq:A1-orientation-reciprocity}
\end{equation}
The multiplier is independent of the central isogeny, vertex-color lattice,
and elliptic embedding phase, although the generating functions differ.
\end{theorem}

\begin{proof}
The classification theorem identifies every local coefficient with the
Frobenius trace of the displayed interpolation.  The saturated identity is
\eqref{eq:A1-saturated-reciprocity}; the type-preserving identity is
Proposition~\ref{prop:SL2-parity-reciprocity}.  Finite central isogenies do
not change the rational root cocharacter space, and
Proposition~\ref{prop:U2-center-test} deletes the connected-center direction.
This proves the formula in every classified case.
\end{proof}

\begin{corollary}[The Kummer correction is structural]
\label{cor:A1-structural-Kummer}
The line $\mathfrak o_{G,T}$ is the sign summand of the quadratic root
cameral cover.  On a chart $v^2=u$ it is the Kummer local system $\cK_u$,
and every allowed positive root shell contains the factor
$(\Lell-\mathfrak o_{G,T})\Lell^{r-1}$ before summation.  In the saturated
$PGL_2$ case this factor is represented by $\mathscr Q_r$; in all other cases
it denotes the orbit-summed interpolation.  If the cameral cover is
nontrivial, the line is not Tate.
\end{corollary}

\begin{proof}
The deck involution acts by $-1$ on the rank-one root cocharacter lattice, so
its determinant is the sign summand.  The represented assertion is
Theorem~\ref{thm:finite-residual-filtration}; the general trace statement is
Lemma~\ref{lem:A1-shell-transfer}.  Nontrivial geometric monodromy excludes
a constant Tate twist, as in Proposition~\ref{prop:strict-non-Tate-Kummer}.
\end{proof}


\section{Literature comparison, limitations, and outlook}
\label{sec:literature-outlook}

\subsection{Comparison and scope of claim}

The split and unramified-elliptic numerical formulas are classical and agree
with the rank-one calculations of Kottwitz \cite{KottwitzOrbital}; affine
Springer geometry in the unramified setting goes much further
\cite{GKM}.  Hales proves that broad classes of orbital integrals have
virtual motivic representatives, and transfer principles compare their
specializations \cite{HalesMotivic,CluckersHalesLoeser}.  Chen's fundamental
domains and truncations relate affine-Springer point counts to ordinary and
weighted orbital integrals
\cite{ChenFundamentalDomain,ChenWeighted,ChenLocalBehavior}.
Kummer local systems already occur in the orbital-integral setting of
Cunningham--Hales \cite{CunninghamHalesGood}.

The theorem proved here is narrower.  It is the componentwise identification
of the compact positive-loop orbit sheaf in the universal quadratic $PGL_2$
family with the single smooth quotient scheme $\mathscr Q_r$ over the whole
parameter base, together with its enhanced determinant-line descent.  The
termwise equality of represented shells, Haar masses, and Frobenius traces is
stronger than a numerical or virtual interpolation, but it is proved only for
this family.  The absolute-$A_1$ theorem makes a different contribution: it
shows that the same root-orientation line survives the $SL_2$ change of shell
lattice and that the $U_2$ center test forces deletion of
$X_*(Z(G)^\circ)$.  Its general coefficients are local orbit-summed
interpolations.  These statements specify the scope of the paper; they are
not accompanied by a claim that no equivalent formulation exists elsewhere.

Stanley reciprocity and cellular Verdier duality are established tools
\cite{Stanley,Curry,BresslerLunts}.  The relation between affine-Springer
point counting and orbital integrals, the $SL_2$ vertex-color parity, and the
exterior/Poincar\'e description of torus cohomology are likewise used as
standard input.  Their role here is to identify and test the height-graded
shadow of the represented orbital family.  The paper does not claim a new
general cellular-duality or torus-duality formalism.

\subsection{Exact limitations}

Theorem A concerns the equal-characteristic loop family over $\Uk$.  It
represents each compact $L^+\mathbb T$ orbit, but it does not represent the
whole affine-Springer locus $X_m$, the quotient stack $[T\backslash X_m]$, or
closures and incidence relations among all its strata.  The bookkeeping
union $\mathscr S_m^{\mathrm{orb}}$ contains one compact fundamental shell per
local double coset and is not a global decomposition theorem.

Theorem B is complete only for Definitions~\ref{def:A1-pair} and
\ref{def:standard-A1-family}.  It excludes ramified tori, anisotropic inner
forms, nonspherical test functions, nonstandard depth/fixed-set relations,
and relative rank-one groups of non-$A_1$ absolute type such as quasi-split
$U_3$.  Central neutrality is essential.  Outside the saturated $PGL_2$
case, no relative fppf orbit scheme or derived orbital push-forward has been
constructed.

Finally, Theorem~\ref{thm:PGL2-genuine-small-skeletonization} has a discrete
finite target.  It is object-level but contains no cross-shell incidence.
The two-dimensional cone in
Theorem~\ref{thm:PGL2-narrow-skeletonization} appears only after
cohomological associated graded, de-Tating, and passage to a rational
$K_0$ character.  There is no infinite cellular object, completed Rees
module, or convolution theorem hidden in that statement.

\subsection{Outlook: the remaining bridge}

The next geometric problem is deliberately narrower than a general tropical
realization.  For the same quadratic $PGL_2$ family, one should construct
actual transition morphisms among the represented shells whose incidence
target is the height cone
\[
 C=\{(m,s)\in\mathbf N^2:0\le s\le m\},
 \qquad s=m-r+1.
\]
The desired refinement would upgrade the present discrete functors to an
incidence-bearing pair and prove, before taking $K_0$,
\[
 \Dcell\Sk_!\simeq\Sk_*\D.
\]
The split tame product charts
$\mathscr Q_r\times_U\widetilde U\simeq\mathbb G_m\times\mathbb A^{r-1}$
give local vertical duality, but they do not yet provide the horizontal
transition maps.  Constructing those maps is the first unresolved step.

In this limited sense, the present result isolates a rank-one geometric
fragment of a possible correspondence-level tropical transport theory.

\section{Conclusion}

The paper proves two results of different strength.  First, the universal
quadratic $PGL_2$ family has a componentwise geometric model:
$\mathscr Q_r$ represents the compact positive-loop orbit of
$[\mathcal A_r]$, and its finite-field points give exactly the Haar mass of
the corresponding spherical double coset.  The residual vector-torsor tower
identifies every compact shell coefficient with the torus Frobenius bimodule
and propagates the non-Tate determinant line
$\cK=\det(X_*(\mathbb T)\otimes\Qell)$ through all depths.  The finite
labelled-shell map supplies a genuine, if discrete, Verdier partner.

Second, classification and orbit-summed measure transfer extend the
reciprocity multiplier to every central-neutral standard unramified
absolute-$A_1$ family.  The $SL_2$ parity test changes the shell lattice while
leaving the multiplier fixed; the $U_2$ center test forces the formula
\[
 \mathfrak o_{G,T}
 =\det\!\left(X_*(T/Z(G)^\circ)\otimes\Qell\right).
\]
Thus the Kummer line in the represented $PGL_2$ family is the rank-one
realization of root orientation, not an accidental numerator in one rational
function.

\section*{Acknowledgments}

This mathematical paper was generated 100\% by artificial intelligence.  I
first acknowledge the textual data produced by humanity as a whole, which
provides the linguistic and cultural record from which mathematical expression
can be learned.  I then thank the accumulated scholarly achievements of
mathematicians throughout history.  I also thank the researchers whose work has
advanced and refined neural networks through successive generations.  Finally,
I thank the OpenAI team for the training and development of the models used in
the preparation of this paper.

\begin{thebibliography}{99}

\bibitem{BresslerLunts}
P.~Bressler and V.~A. Lunts,
\emph{Intersection cohomology on nonrational polytopes},
Compos. Math. \textbf{135} (2003), 245--278;
\url{https://arxiv.org/abs/math/0002006}.

\bibitem{ChenFundamentalDomain}
Z.~Chen,
\emph{On the fundamental domain of affine Springer fibers},
Math. Z. \textbf{286} (2017), 1255--1306;
\url{https://arxiv.org/abs/1303.4630}.

\bibitem{ChenWeighted}
Z.~Chen,
\emph{Truncated affine Springer fibers and Arthur's weighted orbital
integrals},
J. Inst. Math. Jussieu \textbf{22} (2023), 1757--1818;
\url{https://arxiv.org/abs/1701.03202}.

\bibitem{ChenLocalBehavior}
Z.~Chen,
\emph{On the local behavior of weighted orbital integrals and the affine
Springer fibers},
preprint, arXiv:2103.15091;
\url{https://arxiv.org/abs/2103.15091}.

\bibitem{CluckersHalesLoeser}
R.~Cluckers, T.~Hales, and F.~Loeser,
\emph{Transfer principle for the fundamental lemma},
in \emph{On the stabilization of the trace formula}, Stab. Trace Formula,
Shimura Var. Arith. Appl., vol.~1, Int. Press, 2011, 309--347;
\url{https://arxiv.org/abs/0712.0708}.

\bibitem{ConradReductiveGroups}
B.~Conrad, notes by T.~Feng,
\emph{Reductive groups over fields},
course notes, version of 5 December 2020;
\url{https://math.stanford.edu/~conrad/249BW16Page/handouts/AMSalggroups2.pdf}.

\bibitem{CunninghamHalesGood}
C.~Cunningham and T.~C. Hales,
\emph{Good orbital integrals},
Represent. Theory \textbf{8} (2004), 414--457;
\url{https://arxiv.org/abs/math/0311353}.

\bibitem{Curry}
J.~Curry,
\emph{Dualities between cellular sheaves and cosheaves},
J. Pure Appl. Algebra \textbf{222} (2018), 966--993;
\url{https://arxiv.org/abs/1607.06942}.

\bibitem{GKM}
M.~Goresky, R.~Kottwitz, and R.~MacPherson,
\emph{Homology of affine Springer fibers in the unramified case},
Duke Math. J. \textbf{121} (2004), 509--561;
\url{https://arxiv.org/abs/math/0305144}.

\bibitem{HalesMotivic}
T.~C. Hales,
\emph{Orbital integrals are motivic},
Proc. Amer. Math. Soc. \textbf{133} (2005), 1515--1525;
\url{https://arxiv.org/abs/math/0212236}.

\bibitem{KottwitzOrbital}
R.~E. Kottwitz,
\emph{Harmonic analysis on reductive $p$-adic groups and Lie algebras},
in \emph{Harmonic Analysis, the Trace Formula, and Shimura Varieties},
Clay Math. Proc. \textbf{4}, Amer. Math. Soc., 2005, 393--522.

\bibitem{LaszloOlsson}
Y.~Laszlo and M.~Olsson,
\emph{The six operations for sheaves on Artin stacks II: adic coefficients},
Publ. Math. Inst. Hautes \'Etudes Sci. \textbf{107} (2008), 169--210.

\bibitem{LiuZhengEnhanced}
Y.~Liu and W.~Zheng,
\emph{Enhanced six operations and base change theorem for higher Artin
stacks}, arXiv:1211.5948;
\url{https://arxiv.org/abs/1211.5948}.

\bibitem{MilneAlgebraicGroups}
J.~S. Milne,
\emph{Algebraic Groups: The Theory of Group Schemes of Finite Type over a
Field},
Cambridge Studies in Advanced Mathematics 170,
Cambridge University Press, 2017;
\url{https://www.jmilne.org/math/Books/iAG2017.pdf}.

\bibitem{PappasRapoport}
G.~Pappas and M.~Rapoport,
\emph{Twisted loop groups and their affine flag varieties},
Adv. Math. \textbf{219} (2008), 118--198;
\url{https://arxiv.org/abs/math/0607130}.

\bibitem{Stanley}
R.~P. Stanley,
\emph{Combinatorial reciprocity theorems},
Adv. Math. \textbf{14} (1974), 194--253.

\bibitem{ZhuAffineGrassmannians}
X.~Zhu,
\emph{An introduction to affine Grassmannians and the geometric Satake
equivalence},
lecture notes, arXiv:1603.05593;
\url{https://arxiv.org/abs/1603.05593}.

\end{thebibliography}

\end{document}
